\quellseite{406}

\part{Das erste Buch Samuels}

\chapter{Das 1. Kapitel}

\section{Vorrede}

Heute fangen wir an das erste Buch Samuels. In diesem Buch aber werden fürnehmlich gemeldet Exempel, wie Gott auch an gotterleuchteten und großen Leuten die Sünde gestraft, und dagegen, wie er auch an gemeinen und geringen Menschen seine Gnade bewiesen hat. Diese Exempel sollen wir zu täglicher Warnung und Trost fleißig anhören und betrachten. Denn hat Gott große, hochgeachtete Leute gestrafet; fürwahr, er wird unser auch nicht verschonen, da wir sündigen. Also hat Gott unansehnliche Leute für der Welt gesegnet; um unserer Gebrechlichkeit willen wird er unser mit seinem Segen traun auch nicht vergessen, da wir in wahrer Buße leben. Derwegen wollen wir uns alle Augenblicke für Sünden hüten, daß wir nicht straffschuldig werden, und wollen dagegen alle Tage christlich leben, daß uns Gott auch segne und behüte. O HErr JEsu Christ, das hilf uns allen, Amen. Darauf höret ferner an das 1. Kapitel.

\section{Samuel wird geboren}
\quellseite{407}

\section{Summarium}

Ist eine Historia, wie Hanna, als eine gottesfürchtige Frau, Gott um einen Sohn bittet, und Gott sie ihrer Bitte gewähret, V. 20. Auf daß wir erstlich lernen sollen, daß die Kinder eine Gabe Gottes sind, wie der 127. Psalm V. 4. auch lehret. Danach, daß Gott Kinder desto lieber geben will, wo man sie Gott zu Ehren und den Leuten zu Nutz begehrt aufzuziehen; wie denn der heiligen Hanna Herz und Gedanken gestanden, V. 11.

Solchem Fürnehmen der lieben Hanna sollen die Mütter auch folgen, ihre Kinder zu Gottesfurcht und aller Zucht halten. Denn dasselbe heißt nach unserer Sprache, dem HErrn die Kinder geben, nicht in die Klöster stoßen, wie im Papstthum die Gewohnheit war.

\section{Votum}

Helfe GOtt, der Vater aller Gnaden, daß alle Mütter solcher heilsamer Vermahnung treulich folgen, und regiere alle lieben Kinder, daß sie sich ziehen lassen und alle wohl gerathen, GOtt zu Lob und ihrem Nächsten zu Nutz, Amen.

\chapter{Das 2. Kapitel}

\section{Vorrede}

Wir sind Pflanzen Gottes, ihm zu Lob und Preis, sagt Jesaias im 61. Cap. V. 3. Wer aber Gott ehret, den will Gott wieder ehren. Ueber das alle Tage, alle Stunden, ja alle Augenblicke allerlei Gutes an Leib und Seele erzeigt uns Gott so reichlich, daß wir mit dem Patriarchen Jakob bekennen müssen: Wir sind viel zu geringe und unwürdig gegen Gottes große Barmherzigkeit, 1 Mos. 32, 10. Derwegen wir Gott für seine Wohlthaten danken sollen mit Herz und Munde, und auch mit der That, also daß ein jeder in seinem Beruf, als in der Kinderzucht und dergleichen sich christlich verhalte, das Böse lasse, das Gute thue, das Kreuz trage und das Unrecht überwinde, und alles im Glauben und Geduld. Dies erinnert auch das 2. Kapitel.

\section{Hannas Lobgesang. Auferziehung Samuels. Bosheit der Söhne Eli}

a Alte) das Feste, Gewisse, Ehrliche, wie man spricht: Gewohnheit, alt Herkommen, Landessitten und Weise; denn darauf trotzen die Leute, und sagen: Ei sicher, die alte Weise die beste, unsere Vorfahren sind auch nicht Narren gewesen. Und dennoch also wider Gottes Wort, als müßte es nicht anders noch newern.

b Das ist, sie müssen ums Brod dienen.
\quellseite{408}

\section{Summarium}

Ist ein Danklied, in welchem die heilige Frau Hanna Gott lobet, daß er allein der rechte HErr und Gott sei, der alles nach seinem göttlichen Willen könne machen, V. 3. Darum vermahnet sie, V. 6., daß man allein auf solchen Gott trauen, denselben allein fürchten und seine Hülfe allein suchen, und derselben auch mit Geduld erwarten soll. Die es aber nicht thun, sondern ihren Trost auf Geld, Gut und Gewalt setzen, wie die Welt pflegt, die müssen zu Boden gehen.

Am Ende gedenkt sie, wie Gott die Welt richten werde, V. 11., das ist, das Evangelium in aller Welt predigen lassen, und seinem Gesalbten, das ist, unserm HErrn Jesu Christo, Macht über alles geben. Mit solchem Danklied stimmet überein der Gesang der heiligen Jungfrau Maria, und ist zu beiden Theilen einerlei Meinung und schier auch einerlei Wort.

(Leibrock) das waren priesterliche Kleider, davon 2 Mos. 28, 6.

d Si Deus offenditur, et non ipse per se remiserit, non est aliquis alius vel superior, qui inter ipsum offensum et offensorem mediare possit, sicut inter homines. D. h. wenn Gott beleidiget wird, und nicht selbst durch sich selbst vergiebt, so ist kein Anderer oder Höherer, der zwischen ihm, dem Beleidigten und dem Beleidiger vermitteln könnte, wie unter Menschen.
\quellseite{409}

e (Lückest du) gleichwie Apost. 9, 5. St. Paulus wider den Stachel löcket, das ist frech und muthwillig.

f (Altar) Non auferam quidem de altari meo, sed diu non vivent. Quia ubi adoleverint, morientur.

\section{Summarium}

Ist eine Historia von dem Muthwillen und Unzucht der Söhne Eli, des Hohenpriesters, welche Gott also straft, daß, wie folgen wird, nicht allein sie wegen auf einen Tag umkommen, V. 34.; sondern Gott dräuet auch, daß ihr ganz Geschlecht soll solcher Sünde entgelten, daß wenn sie zu Mannesalter kommen, sterben sollen, und also ihrer Nachkommen keiner zum vollkommenen Alter kommen soll, V. 35. Auf daß die Eltern ja lernen ihre Kinder mit Fleiß in Gottesfurcht aufziehen, und ihnen ihren Muthwillen nicht gestatten, weil es so hart und greulich soll gestraft werden, wo sie hierin lässig sind.

\section{Votum}

Regiere uns der treue GOtt, daß wir ihm für seine Wohlthaten dankbar sein, und segne die Kinderzucht, auf daß sie bei der lieben Jugend wohl gerathe, damit sie in Verstand zunehme und in Tugenden lebe, Gott zu Ehren, ihnen selbst und andern Leuten zu Nutz, Amen.

\chapter{Das 3. Kapitel}

\section{Vorrede}

Salomon in seinen Sprüchen im 29. Kapitel V. 15. sagt: Ruthe und Strafe gibt Weisheit; aber ein Knabe ihm selbst gelassen, schändet seine Mutter. Züchtige deinen Sohn, so wird er dich ergötzen, und wird deiner Seele sanft thun. Ferner Sirach im 26. Kapitel V. 13. spricht: Ist deine Tochter nicht schamhaftig, so halt sie hart, auf daß sie nicht ihren Muthwillen treibe, wenn sie so frei ist. Wenn du merkest, daß sie frech um sich stehet, so siehe wohl darauf; wo nicht, und sie thut darüber wider dich, so laß dich es auch nicht wundern. Solche Kinderstrafe aber sollen die Eltern also mäßigen, daß die Kinder selbst merken und spüren mögen, daß ihre lieben Eltern sie nicht aus Haß und Feindschaft gegen ihre Personen, sondern aus lauter Liebe strafen, ihre Sünden anfeuren und sie gern mit ihrer Züchtigung davor verwahren wollten, daß sie der irdische oder geistliche Scharfrichter (und alsdann der höllische Henker und Peiniger) nicht strafen dürfen. Dieses erinnert auch das 3. Kapitel.

\section{Samuel wird berufen und ansehnlich gemacht}

a Theuer: es waren nicht Prediger noch Pfarrherrn genug, die Bibel ist da unter der Bank gelegen, hat niemand studiret; Samuel ist kommen, und hats wieder hervor gebracht.
\quellseite{410}
ᵇ Id est, Deus coepit apparere copioso verbo sub Samuele.

\section{Summarium}

Ist eine Historia, wie Gott sich dem Knaben Samuel offenbaret, V. 11. Und zum andernmal dem Hohenpriester Eli durch ihn dräuet, er wollt seiner Söhne Sünden strafen, und er, der Vater selbst, soll auch darum gestraft werden, daß er gewußt hat, daß seine Söhne sich so schändlich halten, und dennoch nicht einmal dazu sauer gesehen, V. 12. Solches sollen die Eltern mit Fleiß merken. Eli hatte ja seine Söhne gestraft, wie oben im andern Capitel V. 24. stehet; aber Gott will sich nicht daran genügen lassen, wenn man allein mit Worten will strafen, sondern er will, daß man mit Ernst dazu thun und mit harter Strafe der Sünden wehren soll, und dieselben in keinem Weg an den Kindern dulden. Das that Eli nicht, und sagte seinen Söhnen ein Wort, zwei, und ließ es dabei bleiben, brauchte keinen sondern Ernst wider sie, ließ sie ihm zu lieb sein. Das war unrecht, darum mußten Vater und Söhne ihre Strafe darum leiden.

\section{Votum}

HElfe GOtt, der seinen Willen will erfüllet haben, daß die Eltern in der Kinderzucht gebührlichen Glimpf und nöthigen Ernst gebrauchen, und regiere die Kinder, daß sie sich ziehen lassen und gehorsam sein und bleiben allzeit, Amen.
\quellseite{411}

\chapter{Das 4. Kapitel}

\section{Vorrede}

Der meiste Theil der Menschen gehet dahin, und bedenket nicht einmal, was Sünde sei und wie großen Schaden sie mit sich bringet. Die Jugend achtets nicht aus Unverstand und Nachlässigkeit, so wollen es die Alten nicht wissen noch erkennen, vor großer Klugheit und Sicherheit. Derwegen erzählet die heil. Schrift die schrecklichen Sündenstrafen, in viel und mancherlei Exempeln, jedermann zur Warnung. So läßt uns solche anhören, daß ein jeder vor Gottes Zorn erschrecke, Buße thue, die Sünde meide, und dagegen christlich lebe, daß Gott nicht auch also wegen der Sünde uns züchtigen dürfe, wie er andere Land und Leute derwegen grausam gestrafet hat. Darauf höret auch an das 4. Kapitel.

\section[Israels Niederlage. Entführung der Bundeslade. Todesfall Eli und]{Israels Niederlage. Entführung der Bundeslade. Todesfall Eli und seiner Söhne}

a Eben(Ezer) heißt Helfstein. f. unten Kap. 7, 12.
\quellseite{412}

\section{Summarium}

In eine Historia, in welcher man sehet, wie Gott endlich die Sünde hart und zugleich strafet. Und wird oben gemeldet, wie die Leute das Opfer des HErrn lästerten; item, wie Gottes Wort seltsam und wenig Weissagung gewesen sey. Die Sünden nun, so hernach gefolget sind, strafet Gott also, daß die Philister vier und dreißig tausend Mann erschlagen, V. 10., und die Lade des Bundes auch wegnehmen. Die zween gottlosen Söhne des Priesters Eli kommen auch um, V. 11. Der Vater Eli, da er solches höret, erschrickt also, daß er zurückfällt und den Hals bricht, V. 18. Solches sind eitel Strafen der Sünde. Sonderlich aber ist zu merken, weil die Kinder Israel nicht Buße thaten, half sie es nichts, daß sie die Lade des Bundes bei ihnen hatten. Eben wie es uns wider den Türken auch nicht hilft, daß wir Christen heißen, getauft sind, Gottes Wort und das Sacrament haben, sintemal wir unbußfertig sind, und von unserm gottlosen Wesen nicht ablassen.

\section{Votum}

GOtt erbarme dich unser, lindere die wohlverdienten Sündenstrafen; verleihe auch, daß wir in seinem Gehorsam leben, und führe uns aus dem Sündenreich dieser Welt in das Gnadenreich der Seligkeit, um JEsu Christi unsers HErrn willen, Amen.

\chapter{Das 5. Kapitel}

\section{Vorrede}

Was den Menschen unmöglich und schwer ist, das alles ist Gott möglich und leicht, und Gott kann in aller Noth erretten alle, die im wahren Glauben zu ihm treten und um Beistand ihn ansprechen. Dergleichen in schweren Sachen und großer Gefahr niemand zagen oder verzweifeln soll; sondern jedermann habe Zuflucht zu Gott, und suche durchs Gebet bei ihm Hülfe, er wird die Seinen wider ihrer Feinde sichtbare und unsichtbare List und Gewalt wohl behüten und erhalten. Diesen Trost lehret Gott selbst Jes. 43, 2. und sagt: So du durch Wasser gehest, will ich bei dir sein, daß dich die Ströme nicht sollen ersäufen, und so du ins Feuer gehest, sollst du nicht brennen, und die Flamme soll dich nicht anzünden; denn ich bin der HErr, dein Gott, der Heilige in Israel, dein Heiland. Ferner höret an das 5. Kapitel.

\section{Plage der Philister wegen der entführten Bundeslade}

a Wenn Menschen und Vernunft obliegt, so gehet Gottes Wort und alle Ehre dahin; da fället das rechte Priesteramt zurück und stirbet, und das alles aus Gottes Zorn. Aber die so gewinnen, haben darnach keine Ruhe im Gewissen. Denn wo Gottes Wort nicht recht gehet, thut es dem Gewissen alles Leid an, wie hie die Lade Gottes den Philistern, so lange, bis sie zuletzt ihre Schande erkennen müssen, daß sie Gottes Wort verachtet haben und mit Ehren wieder zurecht bringen müssen. Das ist hie bedeutet mit den äußeren Zeichen und Mäusen, welches nichts ist denn die heimliche Plage der Gewissen, die zuletzt offenbar wird durch Gottes Wort, wie St. Paulus sagt 2 Tim. 3, 9.: Ihre Thorheit wird offenbar werden jedermann.
\quellseite{413}

\section{Summarium}

Dies ist eine schöne Historia, wie es den Philistern geraten sei, daß sie die Lade des Bundes genommen haben, nämlich, daß sie Gott also mit Krankheiten und Sterben strafet, daß sie so froh werden, da sie ihrer wieder los wurden, V. 6. Solche Historia reimet sich fein mit den Ketzern, so Gottes Wort den Christen nehmen und zu eigener Ehre, um des Bauches willen, mißbrauchen. Aber endlich strafet sie Gott mit heimlichen Krankheiten, das ist, mit bösem Gewissen und Verzweiflung; das soll ihr Lohn sein.

\section{Votum}

Der barmherzige, starke GOtt steure allen Rotten und Secten, die sein Wort verkehren und seine Kirche zerstören; und behüte uns für Land-, Stadt- und Kirchen-Krieg und Empörungen, um seines Namens Ehre und der aufwachßenden Jugend Wohlfahrt willen, Amen.

\chapter{Das 6. Capitel}

\section{Vorrede}

Der wahre Gottesdienst ist, recht erkennen aus Gottes Wort das Wesen und den Willen Gottes. Das Wesen also, daß nur ein Gott ist, aber drei verschiedene Personen sind. Den Willen aber, daß es der ist, daß wir an den Sohn glauben sollen, und durch den Glauben an ihn erlangen Vergebung der Sünden, Gerechtigkeit, den heil. Geist und ewiges Leben, und diesen wahren Gott, der also erkannt wird, lieben und fürchten, auf ihn hoffen, ihm gehorsam sein im Glück und Unglück, ihn anrufen, ihm danken, seinen Namen predigen, Gottes Wort gern und fleißig hören, das Sacrament oft gebrauchen, den Eltern, der Obrigkeit, den Zuchtmeistern, Herren und Frauen Gehorsam leisten, barmherzig, keusch, züchtig, auch mäßig sein, gerecht und wohlthätig gegen jedermann, vornehmlich aber gegen Witwen und Waisen, wahrhaftig und aufrichtig, und sich unbefleckt von der Welt erhalten. Denn der wahre Gottesdienst begreift in sich alles, was in Gottes Wort, im Gesetz und Evangelio, geboten und befohlen ist. Und diesen Gottesdienst sollen wir alle in gebührlicher Ehre halten, auch christlich verrichten. Als uns denn dieß erinnert das 6. Kapitel.

\section{Die Bundeslade wird den Israeliten wieder zugestellet}
\quellseite{414}

\section{Summarium}

Ist ein fein Exempel, daß die Plage bei den Philistern nicht eher aufhöret, bis sie durch ein Schuldopfer ein öffentlich Bekenntniß über sich geben, daß sie haben Unrecht gethan und sich an dem Gott Israel versündiget, V. 4. Denn soll die Strafe aufhören, so muß man die Sünde erkennen und davon ablassen.

Das aber fünfzig tausend und siebenzig Mann zu BethSemes sterben, darum, daß sie die Lade des HErrn gesehen hatten, V. 10., solches ist nicht darum geschehen, daß sie die Lade schlecht haben angesehen; sondern die Heiden haben ein Frohlocken darüber gehabt, und gedacht, was ohnmächtigen Gott haben doch die Juden, daß er sich und sie nicht rettet? Wir haben das größte Stück von ihrem Gottesdienst und Heiligthum. So es ein rechter Gottesdienst wäre, oder ihr Gott etwas wäre und vermöchte, so hätten sie uns abgeleget. Aber da sehen wir, Gott und das Volk ist eins wie das andere. Werden also einen sondern Stolz gefasset, Gott gespottet und gehöhnet haben, wie solches der Teufel sehr wohl kann. Da läßt sich Gott sehen, mit seinem Zorn, daß ihr fünfzig tausend und siebenzig auf einmal dahin fallen und sterben, und die andern des Spottens aufhören, und wissen nicht, wo sie vor Angst bleiben sollen. Und zum Exempel, daß wir mit Gott, seinem Wort und dem rechten Gottesdienst nicht sollen leichtfertig umgehen, sondern gottesfürchtig sein und in allen Ehren haben.

\section{Votum}

GOtt, der heilige Geist, erleuchte, lehre und führe uns, daß uns der höllische Geist in Verachtung GOttes Worts und der Gottesdienste nicht verführe; und helfe, daß wir in täglicher, christlicher Buße leben, damit uns GOtt unserer Sünde wegen nicht strafen dürfe, Amen.

\chapter{Das 7. Kapitel}

\section{Vorrede}

Rechte Lehrer und gehorsame Zuhörer sind einer jeden Stadt schönster Schmuck und Zierde; und eine solche Gemeinde ist Gottes Paradies, der heiligen Engel Lusthaus und eine Schule und Zuchtspiegel der Jugend. Derwegen danken wir Gott, daß wir sein Wort rein und lauter haben, und daß dennoch dessen auch Hörer und Thäter seien. Der helfe mit Gnaden ferner; was unsere Prediger mit Worten treulich lehren,
\quellseite{415}
daß wir solches mit guten Werken fleißig verbringen, und also leben, daß Gott geprieset, die heiligen Engel erfreuet, die Jugend gebessert und wir von Gott gesegnet werden, und daß wir alle in der Himmelsburg der heiligen Engel und im Lande der Lebendigen und Auserwählten Gott loben und preisen, immer und ewiglich, Amen. Und darauf höret an das 7. Capitel.

\section[Die bußfertigen Israeliten sagen auf Samuels Gebet wider]{Die bußfertigen Israeliten sagen auf Samuels Gebet wider die Philister}

a Weinete) das ist, die Israeliten im ganzen Reich trugen Leid, und klagten dem HErrn über die Philister.

b Das ist, das noch an der Milch ist.

\section{Summarium}

Ist eine sehr schöne Historia, wie Samuel, als ein frommer, rechtschaffener Prediger, ehrlich das Volk, so unter Gottes Zorn und Strafe war, zur Buße vermahnet, daß sie ihre Sünde erkennen, die Abgötterei abthun und ihr Herz zum HErrn bekehren und ihm allein dienen sollen, V. 3. Und richtet darnach den rechten Gottesdienst wieder an und betet für das Volk.

Wo nun also rechtschaffene Buße und Besserung, reine Lehre und rechter Gottesdienst ist, daselbst will Gott gnädig sein und die Strafe abwenden, wie er denn den Kindern Israel hie hilft. Und hier sehet man sonderlich, wie viel an einem frommen, rechtschaffenen Lehrer, der Gottes Wort rein und treulich prediget, und die Leute zum selben vermahnet, gelegen sei. Darum soll man fromme, gelehrte Kirchendiener in allen Ehren und für das größte Kleinod halten und schön aufheben.

Es wird hie im Text gemeldet, die Kinder Israel schöpften Wasser und gossen es aus vor dem HErrn, V. 6. Solches deutet etliche, daß sie vor dem HErrn geweinet haben über die Sünde, damit sie solche Strafe verdienet hätten. Weil aber der Text auch des Fastens gedenkt, muß es verstanden werden vom Wasser, damit
\quellseite{416}
sie sich und ihre Kleider nach Gewohnheit des Gesetzes gereiniget haben. Das aber ist sehr fein, als die Philister ins Land fallen und jedermann furchtsam und verzagt ist, daß das Volk zu Samuel läuft und spricht: Laß nicht ab für uns zu schreien zu dem HErrn, unserm Gott, daß er uns helfe aus der Philister Hand, V. 8. Solche Zuversicht auf Gottes gnädige Hülfe bewegt Gott, daß er ein greuliches Wetter mit Donner und Blitz ließ über die Philister kommen, daß sie flohen und geschlagen worden, V. 18. Darum sollen wirs auch lernen, daß wir in Nöthen und Angst fleißig beten.

\section{Votum}

Der treue GOtt erhalte uns sein Wort und dessen treue Prediger, und regiere uns alle, daß wir recht glauben, christlich leben, das Kreuz geduldig tragen, und daß wir alle selig werden durch JEsum Christum, unsern HErrn, Amen.

\chapter{Das 8. Kapitel}

\section{Vorrede}

Wohlerathene Kinder und fromme Obrigkeit sind zwei edle, große und heilsame Gaben Gottes. Nun rumoret aber erschrecklich der Teufel im Haus- und weltlichen Regiment, und hindert gewaltig das Gute in diesen Ständen. Demgegen bedenke ein jeder gar wohl, was sein Beruf und Amt sei, treibe die treulich und fleißig und bete, daß die Kinderzucht und Regierung ferner glückselig sei, und zu Gottes Ehren und dem menschlichen Geschlecht zu Nutz gerathe und diene. O du Pflanzgarten der lieben Jugend, o du Lustgarten der lieben Obrigkeit, euch segne Gott der Vater, euer warte der treue Gärtner, unser HErr JEsus Christus, Gottes Sohn, euch begegne Gott der heilige Geist, und behüte euch vor giftigen Raupen und allem anderen Ungeziefer, und verleihe, daß ihr wachset, grünet, blühet und Früchte traget, die Gott wohl anriechen und Land und Leuten dienstlich sein, bis zum ewigen Leben, Amen. Darauf höret ferner an das 8. Kapitel.

\section{Israel begehret einen König. Dessen Recht wird angezeigt}
\quellseite{417}

\section{Summarium}

Ist erstlich ein Exempel, wie der fromme, heilige Mann Samuel ungerathene und gottlose Kinder hat, V. 3. Auf daß man sehe, wie auch heilige Leute bisweilen solch Creuz haben tragen müssen. Darum bedarfs nicht allein fleißiger Zucht, sondern auch ernstliches Beten gegen Gott, daß er fromme Kinder geben wolle.

Zum andern ist hie eine Historia, wie die Juden einen König begehren, allein daß sie in solchem dem Heiden gleich seien, V. 5. Solches zeucht Gott für eine Unehre an; denn sie sollten sich mehr ihres Gottes im Himmel, denn eines sterblichen Königs, getröstet haben, V. 7. Es zeiget auch die heil. Schrift hie an, wie die Unterthanen ihrer Obrigkeit schuldig sind zu dienen, nach ihrem besten Vermögen, mit Leib und Gut, auf daß Friede und gut Regiment möge erhalten werden. Aber jedoch soll solches Regiment keine Tyrannei sein, daß man die Leute übel halten, zu viel schätzen und dringen wollte. Denn daß der HErr spricht: Eure Töchter wird er nehmen, und eure Aecker, eure Weinberge und Oelberge, von eurer Saat und Weinberg wird er den Zehenten nehmen. Mit solchen Worten räumet er der Obrigkeit nicht ein, den Unterthanen zu nehmen, was sie gelüstet, sondern will damit anzeigen, daß die Juden solcher Last hätten können gerathen, wo sie an Gottes Schutz sich genügen und keinen König hätten haben wollen. Wo aber ein König ist, da muß jedermann mit dem Seinen zur Unterhaltung des Regiments helfen. Und doch soll des Königs oder Obrigkeit halben die Gleichmäßigkeit bleiben, daß er ein König und nicht ein Tyrann sei. Könige oder herrschafften Obrigkeit lassen den Unterthanen das Ihre, und nehmen ein Leidliches, auf daß es die Leute ertragen und sich und die Ihren nähren mögen. Die Tyrannen aber nehmen es alles und lassen nichts; das gefällt Gott übel und strafets greulich mit Aufruhr und sonst, wie die Exempel weisen.

\section{Votum}

Helffe der treue GOtt, daß sich alle Kinderlein recht ziehen lassen und als Pflänzlein GOttes wohl und nützlich gerathen; und stehe auch unserer lieben Obrigkeit ferner mit Gnaden bei, daß sie christlich und seliglich regiere, in Christo Jesu, unserm HErrn, Amen.

\chapter{Das 9. Capitel}

\section{Vorrede}

Von unsern leiblichen Eltern, als Vater und Mutter, bekommen wir Leib und Leben; von unsern Zucht- und Lehrmeistern werden wir in guten Tugenden und Künsten unterweiset; aber die Kirchendiener und Prediger sind Eltern und Lehrer der Seelen. Wie soll man sie halten? Das lehret Gott selber und sagt durch den Erzpropheten Mose, 5 Mos. 12, 19.: Hüte dich, daß du den Leviten nicht verlässest, so lange du auf Erden lebest. Davon unterrichtet auch Sirach im 7. Kap. V. 33. und spricht: Ehre den Priester und gib ihnen ihr Theil, wie dir geboten ist. Eben also redet St. Paulus, da er spricht 1 Tim. 5, 17.: Die Priester oder Aeltesten, die wohl fürstehen, die halte man zweifacher Ehren werth; nämlich einer höflichen Ehrbietung und ihres Lohnes und Unterhaltung. Wohl denen, die nach diesen Geboten thun und handeln. Denn es wird ihnen gewiß gehen nach der Verheißung des HErrn Christi, Matth. 10, 41., da er sagt: Wer einen Propheten aufnimmt in eines Propheten Namen, der wird eines Propheten Lohn empfangen, und ein Becher kaltes Wasser wird nicht unbelohnet bleiben. Dies lehret auch das 9. Kapitel.

\section{Saul, der ernannte König in Israel, ist Samuels Sach}
\quellseite{418}

\section{Summarium}

Ist eine Historia, wie Saul zum Propheten Samuel kommt; und werden hie sonderlich zwo feine Exempel angezeigt. Also erste, daß dies der feinste und beste Beruf ist, wenn man unversehens und unbedacht zu einem Amt gezogen wird, und sich nicht selbst eindringet, wie Saul hie zu einem König gesalbet wird, da er seine Eselin suchet und gar nichts vom Königreich noch andern dergleichen denket.
\quellseite{419}

\section{Summarium}

Das andere Exempel ist, daß Saul nicht zum Propheten gehen will, er habe ihm denn etwas zu geben. Denn es will Gott sonderlich, daß man die ehren und dankbar gegen ihnen sein soll, welche den Leuten im Wort und geistlichen Amt fürstehen und dienen, wie St. Paulus sehr fein vermahnet, Gal. 6, 6.: Wer unterrichtet wird mit dem Wort, der theile mit allerlei Guts dem, der ihn unterrichtet. Und 1 Cor. 9, 11.: So wir euch das Geistliche säen, ists ein groß Ding, ob wir euer Leibliches ernten? item, der HErr hatte befohlen, daß die das Evangelium verkündigen, sollen sich vom Evangelio nähren. Darum ist es eine große Sünde, die doch zu unsern Zeiten sehr gemein ist, daß man die armen Pfarrherrn auf dem Lande und sonst so kärglich und übel hält. Es wird aber sonderlich in diesem Kapitel das Predigt- oder Kirchenamt sehr hoch gerühmet, da der Text sagt, daß vorzeiten in Israel, wenn man ginge, Gott zu fragen, sprach man: Kommet, laßt uns gehen zum Seher oder Propheten, V. 9. Denn was sie von ihren Propheten und Predigern gehöret haben, haben sie anders nicht gehalten, denn als hätte es Gott selbst vom Himmel herunter geredet. Wie denn Christus will, daß man seine Prediger auch also hören soll, da er spricht: Wer euch höret, der höret mich. Euch, spricht er, das ist, die ihr mein Wort den Leuten fürtragt, und Buße und Vergebung der Sünden in meinem Namen verkündigt, und nicht dem Papst und seine Bischöfe, die das Evangelium lästern und verfolgen.

\section{Votum}

Gott regiere alle christlichen Lehrer, daß sie in ihrem Amte treu sein, und helfe, daß die Zuhörer sich auch aller Gebühr verhalten, und zeitlichen und ewigen Segen davon haben, in Christo JEsu, unserm HErrn, Amen.

\chapter{Das 10. Kapitel}

\section{Vorrede}

Es ist einem gemeinen Landrichter oder Schulheiß tröstlich, daß, wenn ihm die Sache zu hoch und über seinen Verstand und Vermögen ist, er sich derhalben, zur Erhaltung gutes Gewissens, bei seinem Amtmann, wie billig, Raths und Hülfe erholen kann. Vielmehr soll es aller hohen und niedrigen Obrigkeit tröstlich sein, weil bei den geschwinden Zeiten und unbändigen Leuten die Sachen viel zu hoch, schwer und groß sind, über allen menschlichen Witz und Vermögen, (wie alle christlichen Regenten und Obrigkeiten gern bekennen,) daß sie derhalben zu Hofe suppliciren und bei der höchsten Gewalt, bei dem HErrn, welches Reichs Amtleute sie sind, sich Raths, Hülfe und Beistand erholen mögen. Und eben das will Gott, der HErr, auch haben, und gefällt ihm von aller Obrigkeit und Menschen herzlich wohl, da man nur getrost Rath und Hülfe bei seinen göttlichen Gnaden sich erholet. So helfe Gott, der heilige Geist, daß alle Obrigkeit des sich christlich verhalte, zu ihrem und der Unterthanen Trost und Nutzen, Amen. Dies lehret auch das 10. Kapitel.

\section{Saul zum Könige gesalbet und vorgestellet}
\quellseite{420}

a Wer ist ihr Vater? das ist, lasse sie weissagen, ist's doch nicht vom Vater angeboren, sondern von GOtt, der ist der rechte Vater.

\section{Summarium}

Ist eine Historia, wie Saul zum Könige von Samuel gesalbet wird. Auf daß aber Saul solchem Beruf glaubet, sagt ihm der Prophet etliche Zeichen, auf daß er an denselbigen lerne, Gottes Wort werde ihm nicht lügen, er werde König werden, V. 2.

Zum andern ist sonderlich zu merken, ehe Gott den Saul ins Regiment setzet, gibt er ihm seinen heil. Geist und macht einen andern Mann aus ihm, V. 9. Denn zu einem Regenten, der Glück haben und den Sachen allenthalben soll recht thun, gehöret ein sonderlicher Segen Gottes. Und will sonderlich vonnöthen sein, daß er gottesfürchtig sei und den heil. Geist habe, der ihn recht leite und führe. Denn sonst sind der Sachen und Händel so viel, so sind wir Menschen so schwach und blöde, daß nicht möglich ist, daß ein Regent allzeit recht thue und alle Händel wohl ausrichte; man thut ihm entweder zu viel oder zu wenig. Darum will es vonnöthen sein, daß die Regenten Gott um seinen heil. Geist bitten und sonderlich in Gottesfurcht leben.

\section{Summarium}

Ist ein fein Exempel, ob Gott wohl dem Propheten Samuel genugsam geoffenbaret hatte, er wollte Saul zum König machen; so nimmt dennoch Samuel solche Erwählung öffentlich und nach ordentlicher Weise für, auf daß niemand sich ob solcher Wahl zu beklagen habe, und Saul öffentlich könnte beweisen.

Daß aber Saul sich verkriecht, weil man wählet, V. 23., ist eine Anzeigung, daß ihm für dem Regiment gegrauet habe, als für einem gefährlichen, mühsamen und schweren Amt, wie es denn in der Wahrheit auch ist, und ihm bald im Anfang begegnet, daß ihn etliche verachten. Aber ein weiser Regent soll sonderlich im Anfang viel
\quellseite{421}
lassen für Ohren gehen, und nichts desto weniger seines Berufs mit Fleiß warten; damit wehret man den bösen Mäulern am besten. Wiederum, wer alle Hunde, so ihn auf den Gassen anbellen, schlagen oder werfen will, der wird ihm viel zu schaffen machen und langsam heimkommen. Man lasse sie bellen und warte ein jeder seines Weges und Berufs, so hören sie von ihnen selbst auf, wie hie mit Saul geschieht.

\section{Votum}

Regiere der treue Gott ferner auch unsere liebe Obrigkeit mit seinem heiligen Geist, daß unter ihnen ihm eine Kirche gesammlet und gemeiner Friede erhalten werde, und bewahre sie mit allen Christgläubigen zum ewigen, freudenreichen Leben, Amen.

\chapter{Das 11. Kapitel}

\section{Vorrede}

St. Paulus Röm. 13, 7. lehret, daß man der Obrigkeit schuldige Ehre thun soll. Diese Ehre bestehet darinnen, daß man der Obrigkeit als Gottes Dienerin alle Ehrbietung und Dankbarkeit erzeige; darnach, daß man ihr auch zum Besten gedenke und das Ehrlichste von ihr rede; denn daß man Gott für sie anrufe und auch ihre Gebrechen zudecke, zum Besten deute und zu gut halte. Bei den Römern und Griechen hielt man daher viel darauf, daß der Obrigkeit ihre gebührliche Ehre auch in äußerlichen Geberden erzeiget würde: und wenn sich einer darinnen versehen und etwa einer Person im Regiment nicht gebührliche Reverenz mit Bücken, Neigen und Weichen erzeiget hätte, den hielt man für so gar leichtfertig, daß ihm auch keine Schande noch Laster zu viel sein würde. Gedenket, was will nun uns Christen gebühren, da wir menschliche Vernunft gebrauchen wollen und wider Gottes Gebot nicht thun sollen. Dies erinnert auch das 11. Kapitel.

\section{Sauls erster Sieg wider die Ammoniter}
\quellseite{422}

\section{Summarium}

Ist eine Historia, in welcher man siehet, wie es dem neuen König Saul im Anfang seines Regiments gegangen ist, daß er verachtet ist gewesen und der meiste Theil den Spott aus ihm getrieben. Wie denn alle Sachen, so Gott fürnimmt, zum ersten nicht unangefochten bleiben, sondern gering, schwach und verachtet anzusehen sind. Aber endlich findet man, daß Glück und Heil dabei ist. Wie denn Saul hie mit einem fürstlichen Zorn zusammen bringet und den Feinden absaget, und also auf ein neues zum König erwählet und bestätiget wird. Die aber, so ihn verachtet und übel nachgeredet hatten, müssen sich ducken, V. 15.

Merke aber mit Fleiß, daß der Text spricht: Der Geist Gottes geriet über Saul und sein Zorn ergrimmete sehr, V. 6. Denn solches ist eine Gabe und Werk des heil. Geistes, wo weltliche Obrigkeit sich mit Ernst darum annimmt, daß den Unterthanen geholfen und die bösen Buben gestraft werden. Darum jedermann solche Regenten für Augen haben, Gott für sie danken und bitten soll, daß Gott Glück und langes Leben ihnen geben und sie für Unfall behüten wolle.

\section{Votum}

Helfe Gott, der zu gebieten und zu verbieten allein Macht hat, daß der Obrigkeit und Unterthanen Herz gleich und einig sei, und daß die Obrigkeit väterlich und die Unterthanen kindlich sich gegen einander erzeigen, und gebe unserer Obrigkeit auch Glück und Sieg wider alle Feinde, Amen.

\chapter{Das 12. Kapitel}

\section{Vorrede}

Aller getreuer Prediger Amt ist, daß sie strafen, das ist, allerlei Sünde wider die erste und andere Tafel anzeigen, Gottes Zorn und Rache darüber verkündigen und zur Besserung vermahnen sollen. Denn der HErr Christus hat befohlen Luc. 24, 47, daß man in seinem Namen predigen solle Buße und Vergebung der Sünden. So kann Gottes Wort ohne Strafen nicht recht fürgetragen werden. Denn Gottes Wort zwei Hauptstücke lehret, nämlich Gesetz und Evangelium; das Gesetz soll strafen, das Evangelium soll trösten. Ueber das, so fordert's aller derjenigen Nothdurft, welche sollen wiedergeboren, gerecht und selig werden, daß sie erstlich durch Strafe des Gesetzes zur Erkenntniß der Sünden kommen, darnach, daß sie durch die Gnadenpredigt des Evangelii wieder getröstet werden. Davon lehret auch das 12. Kapitel.

\section{Wie anfänglich Samuel sein Richteramt übergeben}
\quellseite{423}

a Gott bestätiget den König und zürnet doch, daß sie ihn erwählet hatten. Das ist so viel: sie hätten übel, daß sie ihr Vertrauen von Gott auf einen Menschen und sich selbst setzten, so sie bisher so oft ohne König, durch Gott errettet waren. Dazu, weil ihnen versprochen war, Könige zu haben, läßt es Gott nicht zu, daß sie ihn wählen, sondern er selbst wählet, auf daß bestehe, daß alles, was Gott nicht ansiehet und thut, nichts gelte für Gott.

\section{Summarium}

Ist ein sehr fein Exempel von dem Propheten Samuel, in welchem man siehet, wie ein rechter Prediger oder Kirchendiener sich halten soll: nichts um Geldes und Bauches willen thun, sondern mit allem Fleiß den Leuten ihre Sünde und Untugend fürhalten, sie zur Besserung vermahnen und selbst ihnen gute Exempel geben; auf daß der Ruhm bestehen möge, daß in solchem Amt niemand ihm selbst, sondern Gott und dem Nächsten gedienet habe. Man zeuge jemand wider mich, spricht er, V. 3., ob ich jemandes Ochsen oder Esel genommen, ob ich jemand Gewalt und Unrecht gethan, ob ich jemand unterdrückt, ob ich von jemand Geschenk genommen habe, so will ich's wieder geben. Denn das sind gemeiniglich deren, so in Aemtern sind und sonderbare Verwaltung haben, böse Sünde, daß sie mit Gewalt fahren, andere drücken und sich bestechen lassen. Daher denn das Sprichwort kommt: Es sei kein Aemtlein so gering, es sei Heutens werth.

Ein Prediger aber soll sonderlich für dergleichen Aergerniß sich hüten. Denn sonst verursachet er, daß beides, das Amt und das Wort verachtet und gelästert wird, das ist eine große Sünde.

\section{Summarium}

Ist eine sonderliche Historia, wie Samuel, als ein rechtschaffener Prediger, seinem Volk ihre Sünde anzeigt. Es war aber wider Gott und seine Ehre, daß das Volk einen König wollte haben, und sich mehr eines Königs, als Gottes tröstet. Solche Sünde erkannten sie nicht, ließens ihnen auch nicht leid sein, darum zeiget ihnen Samuel solche Sünde an und sagt V. 17., Gott habe es also gestrafet, daß zu einer ungesegneten Zeit, nämlich in der Weizenernte, Regen, Hagel und Ungewitter sei kommen. Wie denn Gott noch unsere Sünde mit Ungewitter, mit großer Dürre oder Nässe, daß wir der Früchte nicht genießen können, strafet, doch ist solche Strafe weit gnädiger, denn Krieg oder langwierige theure Zeit.

Wiederum thut Samuel auch in dem, wie ein frommer Prediger, daß er sein Volk tröstet, und sie vermahnet, daß sie doch fortan wollen frömmer sein, so wolle er auch gerne mit Beten und Lehren das Beste bei ihnen thun.

\section{Votum}

GOtt, der HErr, bezahle allen christlichen Predigern und Seelsorgern ihre väterliche
\quellseite{424}
Liebe und Treue, und regiere uns, daß wir ihre Straffen annehmen, daraus uns bessern und durch das heilige Evangelium wider die Sünde uns mit Freuden trösten, im Leben und Sterben, Amen.

\chapter{Das 13. Kapitel}

\section{Vorrede}

Der königliche Prophet David sagt im 119. Psalm V. 11.: Herr, ich behalte dein Wort in meinem Herzen, auf daß ich nicht wider dich sündige. Wie nun David Gottes Wort nicht allein in Büchern und im Munde, sondern auch in seinem Herzen, als in einem Kasten behält und durch Gottes Hülfe bewahret, auf daß er nicht wider seinen lieben Gott sündige: also sollen wir auch aus Gott sein und nicht allein Gottes Wort hören, sondern auch in einem feinen, reinen, guten Herzen bewahren und darnach leben, Joh. 8, 47. Luc. 8, 15. Denn wer sich für Sünden hüten will, der kann kein besser Mittel finden, denn daß er Gottes Wort in seinem Herzen verschließe, und sehnlich daneben Gott anrufe, daß er ihn dabei gnädig erhalte, und daß er Gottes Gebot nicht aus den Augen sehe, sondern sein Leben, Wesen, Wandel und Handel als nach dem rechten unbetrüglichen Winkelmaß oder Richtscheid anstelle, daß sein Thun und Lassen allzeit mit demselbigen überein komme. Dies soll auch erinnern das 13. Kapitel.

\section{Sauls Mißtrauen im Zuge wider die Philister}
\quellseite{425}

\section{Summarium}

In eine Historia, wie König Saul sich versündiget, in dem, daß er des Propheten Samuelis nicht erwartet und für sich selbst wider den Befehl Gottes allein opfert, V. 9. Mit solcher Sünde verdienet er, daß Gott ihm und seinem ganzen Geschlecht dräuet, es soll ihr keiner zu königlichen Würden kommen, V. 13. Solches Exempel soll jedermann, insonderheit aber die Obrigkeit, wohl merken und bedenken, wie bald man Glück und Heil verscherzt habe, wenn man etwas wider Gottes Wort und Befehl fürnimmt, aus eigener Andacht. Auf daß man sich desto fleißiger an das Wort halte, und ohne und wider dasselbige seinen sondern Gottesdienst anfange.

\section{Votum}

Regiere und führe uns unser HErr JEsus Christus durch den heiligen Geist, daß wir seinem Wort folgen mit Gehorsam in diesem Leben, daß wir auch seinem Wort und Stimme folgen, mit Freuden in das ewige Leben, und allezeit bei ihm sein, Amen.

\chapter{Das 14. Capitel}

\section{Vorrede}

Geliebte Christen! gleichwie Jonathan, da er die Philister alleine mit seinem Waffenträger angreifen wollte, sich deß tröstete, daß er beschnitten, die Philister aber unbeschnitten wären, und stärket sich also mit dem Sacrament der Beschneidung in seinem Glauben, an Gottes Gnade, Verheißung und Beistand, und hat aus starkem Glauben durch Gottes Hülfe die Philister in die Flucht geschlagen und vertrieben: also sollen wir Christen auch in wahrem Glauben, in Noth und Tod, trösten unserer Taufe, welche anstatt der Beschneidung kommen ist, und sollen uns im Leben und Sterben Gott befehlen, auf seine uns darinnen gethane Verheißung, so werden wir auch siegen über Sünde, Tod, Teufel, Hölle und ewige Verdammniß, und werden die alle vertreiben und selig werden durch Jesum Christum, unsern HErrn. Dies soll man auch betrachten bei dem ersten Theil des 14. Capitels.

\section[Heldenthat Jonathans. Saul führt ihn durch unzeitigen Eifer]{Heldenthat Jonathans. Saul führt ihn durch unzeitigen Eifer in Lebensgefahr}
\quellseite{426}

\section{Summarium}

Ist eine sonderliche schöne Historia, wie Gott das Herz Jonathans erwecket und durch ihn allein die Feinde wunderbarlich dämpfet, V. 14. Denn weil Saul sich versündiget hatte an Gott, ward er solches Siegs nicht werth. Aber sein Sohn war an solcher Sünde nicht schuldig, darum gönnet ihm Gott solchen herrlichen Sieg. Und hier siehet man, wie Gott mit tyrannischen Regenten umgehet und bezahlet sie, wie sie verdienet haben. Die Philister hatten die Juden entblößet, daß sie gar keine Waffen noch Schmied im Lande hatten. Solche Tyrannei strafet Gott damit, daß die Philister ihre eigenen Waffen wider sich selbst brauchen und damit unter einander sich selbst erwürgen müssen.

\section{Votum}

Der allmächtige, barmherzige, starke GOtt stärke heute auch alle christliche Obrigkeit zu glückseliger Regierung, und steure allen Feinden seiner Christenheit, geistlichen und leiblichen, sichtbaren und unsichtbaren, und erhalte uns den leiblichen und geistlichen Seelenfrieden. Amen.

\section{Der andere Theil des 14. Kapitels}

\section{Vorrede}

Geliebte Christen! die Eide und Gelübde, so wider Gott und sein Wort, wider den christlichen Glauben, wider unser Taufgelübniß, wider alle Ehrbarkeit und Billigkeit, oder die da unmöglich sein, die soll man nicht thun noch schwören; und da solches geschehen, soll man sie keineswegs halten; denn wie die Juristen sagen: Eide sollen nicht Bande sein der Ungerechtigkeit. Und der alte Isidorus spricht: In bösem Zusagen halte nicht Glauben und verändere dein Gelübde, wenn es schädlich ist. Und da man einen Eid oder Gelübniß hält, so aus Leichtfertigkeit oder Unverstand gethan ist, und die wider Gott sein und dem Nächsten zu Schaden gedeihen möchten, so macht man die Sünde größer. Und Gott, der alle falsche Eide hasset, strafet solche Sünde an allen unbußfertigen zeitlich und ewig. Derwegen im Eide schwören oder Gelübden ein jeder christliche Bescheidenheit gebrauchen soll. Solches lehret auch der andere Theil des 14. Kapitels.
% Scan page 427: Bible verses/narrative, running heads, folios and bare cross-references only; no non-biblical apparatus retained.
\quellseite{428}

\section{Summarium}

Ist eine Historia, wie Saul als ein unbedächtiger Regent das Volk mit einem Eide verbindet, daß niemand nichts für Nachts essen soll, oder müsse sterben. Ob nun wohl Jonathan wider solchen Eid unwissend that, V. 27., schickts gleichwohl Gott also, daß ihn das Volk errettet, daß er nicht sterben mußte. Zum Zeugniß, daß solcher Eid Gott leid gewesen, und derhalben nicht geleistet hat sollen werden, V. 28. Wir sollen aber dabei lernen, wie gefährlich es sei, und so leichtlich ein Unglück möge erreget werden, wo man mit den Gelübden unweislich und nicht nach dem Wort Gottes sich hält und handelt. Denn eben darum muß die Schuld über Sauls Sohn kommen, auf daß solche unbedachte und ohne Gottes Befehl fürgenommene Geistlichkeit an Saul selbst gestraft werde, der ein Heuchler ist und gedenkt, Gott etwas damit abzugewinnen, so alles Volk den ganzen Tag faste. So doch wir, die da arbeiten sollen, müssen zu essen haben, und Gott nicht das Fasten, sondern das Vertrauen auf seine Güte und Liebe gegen den Nächsten und Besserung des Lebens fordert. Dies sind die rechten Stücke, damit man zum Sieg und Gnade Gottes kommt.

\section{Votum}

Der treue Gott behüte uns für aller Heuchelei, und regiere uns, daß wir uns im Glauben gegen ihn und in der Liebe gegen unsern Nächsten recht verhalten, und errete alle, die unschuldig in Gefahr und Noth stecken, Amen.

\chapter{Das 15. Kapitel}

\section{Vorrede}

St. Petrus spricht in seiner zweiten Epistel im 1. Cap. V. 19.: Ihr thut wohl, daß ihr auf unser festes prophetisches Wort achtet, als auf ein Licht, das da scheinet an einem dunkeln Ort. Hiermit lehret der heilige Apostel, daß die Gottesdienste nichts sind und nichts gelten und nützen ohne Gottes Wort. Dies lehret auch unser HErr Christus, und saget Matth. 15, 9.: Vergeblich dienen sie mir, dieweil sie lehren solche Lehre, die nichts denn Menschen-Gebote sind. Derwegen, wer in seinem Beruf handeln und wandeln will, daß Gott gefalle und ihm und dem Nächsten nützlich sei, der soll sich an Gottes Wort halten und nach dem allein sein Thun und seines Berufes Werk anfangen und treiben. Das lehret auch der erste Theil des 15. Capitels

\section{Saul wegen seines Ungehorsams verworfen}
\quellseite{429}

\section{Summarium}

Dies ist eine Historia, wie der König Saul zum andernmal sich öffentlich versündiget und wider den Befehl Gottes handelt, in dem, daß er der Amalekiter König nicht erwürget und das beste Vieh nicht verbannt, V. 9., wie doch Gott ihm zu thun befohlen hatte, V. 3. Und hilft ihm gar nichts, daß er fürwendet, er habe das beste Vieh behalten zum Opfer dem HErrn, V. 15. Denn der fürnehmste und größte Gottesdienst ist, wie hernach folget, gehorsam sein und thun, was Gott heißet, V. 22. Wer aber solches anstehen lässet und dem Wort ungehorsam ist, er nehme für Gottesdienst, was er wolle, so heißt es und bleibet nur ein Greuel und Sünde für Gott.

\section{Votum}

Der treue GOtt helfe und regiere uns, daß wir uns in all unserm Thun und Lassen sein Wort fürleuchten lassen, auf daß wir wider seinen Willen ja nichts handeln, und durch unser sündlich Gutdünken uns und andern Leuten nicht seinen Zorn und Strafe zuziehen, Amen.

\section{Der andere Theil des 15. Kapitels}

\section{Vorrede}

Eine fürnehme Ursache, um welcher willen Land und Leute zu Boden gehen und die Kirche Gottes Anstoß kriegt, die Regimente auch verderbt werden, ist der Ungehorsam gegen Gottes Wort. Derwegen, wenn treue Kirchendiener und Lehrer allerlei Sünde an ihren Zuhörern sehen, sollen sie, je eher je besser, mit ihnen zur Buße eilen. Und wie sie die Unterthanen ihrer Sünde wegen strafen: also sollen sie zu den Sünden der Obrigkeit auch nicht stille schweigen und sie unterschlagen, sondern sollen ihnen auch Buße predigen, damit beiderseits gelehret und gethan werde, was Gott, das Amt, das Gewissen und eines jeden Pflicht, damit man Gott verbunden ist, erfordert, auf daß man die vorherberichten Strafen abwende und durch den HErrn Christum Gott versöhne, damit er dem geistlichen und weltlichen Stande in Gnaden zu glückseligem Zustande allzeit beiwohne. Davon lehret der andere Theil des 15. Kapitels.
\quellseite{430}

a Man soll Gottes Wort nicht ändern noch bessern; es lässet sich nicht ändern.

\section{Summarium}

Hier sind zwei Exempel sonderlich zu merken. Das erste mit Saul, der thut, wie alle Heuchler pflegen, und will noch Recht haben, seine Sünde schmücken und vertheidigen, V. 20. Dagegen thut Samuel, wie einem rechten Prediger zustehet, schonet Sauls nicht, darum, daß er König war, sondern strafet seine Sünde mit Ernst und sagt: Wer Gott rechtschaffen wolle dienen, der muß seinem Wort und Befehl Gehorsam leisten, V. 22. An selbigem Dienst liegt Gott mehr, denn an allen Opfern; wiederum, wer Gottes Wort ungehorsam sei, er diene Gott, womit er wolle, so sei er um des Ungehorsams willen ein Abgöttischer und Götzendiener. Solchen Spruch sollten Mönche und Nonnen, so sie es werth wären, für sich nehmen, so würden sie befinden, daß ihr Gehorsam in Klöstern der größte Ungehorsam gegen Gott ist, würden derhalben Gelübde, Gehorsam und eines mit dem andern fallen lassen, und sich an Gottes Wort und die zehn Gebote mit rechtem Gehorsam begeben. Da könnten sie alßdenn gewiß wissen, daß sie Gott zu Gefallen und zum Gehorsam lebten, V. 33. Am Ende sehet man, wie Gott den Tyrannen Agag auch bezahlet, wie er verdienet hatte. Und lehret, daß die Feinde der christlichen Kirche endlich ein schrecklich Ende nehmen. Auch sehen wir, was Gott im andern Buch Mose im 17. Kapitel V. 14. wider die Amalekiter hat weissagen lassen; das wird alhier erfüllet. Auf daß wir lernen, wie Gottes Strafen allen Unbußfertigen gewiß folgen, unangesehen, daß sie Gott gleich eine Zeitlang aufheucht.

\section{Votum}

Der barmherzige, gütige GOtt und Vater verleihe, daß Prediger und Zuhörer, Obere und Untere, ihre Aemter treulich wahrnehmen, und daß ein jedes darin GOtt schuldigen Gehorsam leiste; und lindere die wohlverdienten Strafen um Christi willen, Amen.

\chapter{Das 16. Kapitel}

\section{Vorrede}

Man schreibt von einem Wurm, den man Spoliator, das ist, einen Räuber, nennet, daß er auf eine Schlange liegt, die im Schatten liegt, und kriecht ihr erst auf den Schwanz, und fesselt sie fein sanft, zuletzt aber zernaget er ihr das Gehirn und tödtet sie: also ist die Sünde der rechte Räuber und Mörder, sie thut den Menschen wohl sanft, weil er sie begehet nach seines argen Fleisches Lust und Wohlgefallen; aber zuletzt tödtet sie ihn. Denn wie das Quecksilber einen Tam zertreibet und läuft durch alle Glieder des Menschen, und zerreißt sie ihm, wenn es nicht wohl getödtet wird: also verderbet die Sünde den ganzen Menschen innerlich und äußerlich, wenn man die mit Gottes Wort nicht tödtet und ihr widerstrebet. Derwegen sollen wir die vergängliche Lust der Welt fliehen, 1 Pet. 2, 11. Und sollen den alten Adam kreuzigen, tödten und begraben, und sollen nicht durch die Freiheit dem Fleisch Raum geben, sondern im Geist wandeln und die Lüste des Fleisches nicht vollbringen, Röm. 6, 6., Gal. 5, 13. 16., sondern in Heiligkeit und Gerechtigkeit leben. Davon höret auch an das 16. Kapitel.

\section{Davids Beruf und erste Salbung zum Könige}
\quellseite{431}

\section{Summarium}

Ist eine Historia, wie David vom HErrn zum König erwählet und von Samuel gesalbet ist worden, V. 13. Und hie gehet die Historia an, wie die Sünde Sauls sei gestraft worden, aufs erste damit, daß GOtt ihm seinen Geist entzeucht, und lässet ihn treiben von einem bösen Geist, der ihn irre und unruhig machet, daß er hinfort nimmer zufrieden kann sein, V. 14. Auf daß wir lernen die Sünde fliehen und uns dafür hüten, sonsten sie GOtt so greulich strafet; wie wir von Saul hernach weiter hören werden. Durch diese Ursache aber kommet David gen Hof und wird bei dem König und ganzen Volk bekannt.
\quellseite{432}

\section{Votum}

Gott, der heilige Geist, regiere und stärke uns, daß wir den befleckten Rock des Fleisches hassen, und die Werke der Finsterniß ablegen und anziehen die Waffen des Lichts, damit wir nicht auch zeitlich gestraft und ewig verdammt werden, Amen.

\chapter{Das 17. Kapitel}

\section{Vorrede}

Gott ist mächtig, auch in den geringen Menschen große Dinge zu thun. Ja, Gott hat im steten Gebrauch gehabt, die allergrößten Dinge durch unansehnliche Personen auszurichten, mehr, denn durch ansehnliche. Derwegen niemand auf seine Stärke, Macht und Gewalt sich verlassen soll; auch soll niemand in seiner Schwachheit verzagen, wenn er der Dinge sich annimmt, die Gottes Ehre und der Kirchen und der Christenheit Nutz angehen. Denn Gott haffet, was die Welt hoch und groß achtet, und liebet, was für der Welt gering ist, da man nur in seinem Stande demüthig ist, sich und die Verrichtung seines Berufs Gott mit gläubigem Gebete befiehlt, und sich tröstet in wahrem Glauben, daß er durch einen jeden ausrichten könne und wolle, was zu seiner Ehre und andern zu Hülfe, Rettung und Trost dienen soll. Dies erinnert auch der erste Theil des 17. Kapitels.

\section{Davids Streit mit dem Riesen Goliath}
\quellseite{433}

\section{Summarium}

Ist eine Historia, wie David, voll Geistes und Glaubens, mit der Hülfe Gottes den Philister Goliath erwürget, V. 50. Und hie ist im ersten Theil ein Exempel an Goliath wohl zu merken, in welchem uns der Welt Art wird fürgemalet. Denn gleich wie Goliath auf seine Stärke und Größe trotzet, frech und sicher ist, gleich als bedürfe er unsers Herrn Gottes Hülfe gar nicht: also thut die Welt auch, lebt sicher ohne Gottesfurcht, denkt, wenn sie Geld und Gut genug habe, so habe sie alles und bedürfe nichts mehr. Aber wie solcher Trotz und Hochmuth ein Ende nehme, werden wir hernach hören, auf daß jedermann lerne, sich nichts anders rühmen, denn Gottes, seiner Gnade und Barmherzigkeit.

\section{Summarium}

Ist ein fein Exempel, in welchem man siehet, wie Gott seine höchste und gewisseste Hülfe erweisen anschicket, daß es für der Welt ganz und gar kein Ansehen hat. Denn weltlich davon zu reden, wäre mehr Hoffnung dabei gewesen, wenn Saul oder sonst ein großer, starker Mann sich wider den Philister zu streiten hätte brauchen lassen. Aber daß sich David unterstehet, ohne all Schwert und Harnisch mit seiner Schleuder
\quellseite{434}
und Hirtenstab, das scheinet über die Maße liederlich und kindisch. Daher schilt ihn sein ältester Bruder, V. 28. Und der König Saul hat auch geringen Glauben dazu, V. 33. Aber Gott pflegt allweg zu thun, daß er große Wunderwerke durch geringe, verachtete Leute ausrichtet. Erstlich darum, daß die Welt Ursache genug habe, sicher zu sein und Gottes, der seine Sachen schwächlich und thöricht fürnimmt, zu spotten und ihn zu verlachen. Zum andern und fürnehmlich darum, daß man sehe, ja greifen müsse, daß es Gottes Werk und Hülfe sei, und nicht menschlich Vermögen, und also an dem Gegentheil jedermann sehe, daß man mit menschlicher Weisheit, Macht und Vermögen müsse zu Schaden werden, wo Gott nicht dabei ist.

\section{Votum}

Der starke, allmächtige GOtt steure der sichern Welt Trotz und Hochmuth, und stärke seine verachtete Christen, daß sie selige Werkzeuge seiner Kirchen sein und, ihm zu Lob, und Land und Leuten zu Nutz, Freude und Trost, leben, Amen.

\section{Der andere Theil des 17. Capitels}

\section{Vorrede}

Davids Waffen wider Goliath sind ein Stab und Steine in der Hirtentasche; also sollen wir auch den höllischen Goliath, den Teufel, überwinden. Der Stab bedeutet Gottes Wort, die Steine den Glauben, die Hirtentasche das Herz. So sind nun die rechten und kräftigen Mittel, den Teufel zu schlagen und zu überwinden, daß man Gottes Wort und Verheißungen von Herzen glaube und derer wider den Tausendkünstler und alle seine Anläufe, Sturm und Anfechtungen, sich beständig, in Kraft des heil. Geistes, tröste. Derwegen wohl dem, der in Kreuz und Anfechtung wider den Satan also streitet. Denn wie der Apostel Jacobus sagt 1. Kapitel V. 12.: Selig ist der Mann, der die Anfechtung erduldet; denn nachdem er bewähret ist, wird er die Krone des Lebens empfahen, welche Gott verheißen hat denen, die ihn lieb haben. Dies soll auch erinnern der andere Theil des 17. Kapitels.
\quellseite{435}

\section{Summarium}

Ist das herrlichste und beste Stück aus dieser Historia, wie David den Goliath erwürget, nicht durch Leibes Kraft, denn Goliath war mit Stärke dem David und allem Volk weit überlegen, sondern dadurch, daß er sich seines Gottes und seiner Hülfe tröstet. Hierin war David dem Goliath überlegen, der sich allein auf seines Leibes Stärke und Größe verließ, und hatte Gott nicht bei sich. Wie denn David zu ihm sagt: Du kommest zu mir mit Schwert, Spieß und Schild; ich aber komme zu dir im Namen des HErrn Zebaoth, V. 45. Wer ein solch Vertrauen auf Gottes Hülfe und Gnade fassen kann, dem wird der Teufel nicht vermögen etwas zu thun.

Und siehet man sonderlich sein an Goliath, wie es so sehr bald um die Welt mit aller ihrer Stärke, Geld und Gut geschehen ist, wenn Gott nicht dabei ist.

Zu dem sehen wir fürnehmlich in dieser schönen Historia, wie daß David eine Figur Christi unsers HErrn und Heilandes ist, der den rechten Goliath, den Teufel, überwunden, sein Reich zerstöret, die Sünde getilget, den Tod erwürget und die Hölle zerbrochen hat, zu gut, Trost und Seligkeit allen denen, so an ihn glauben, seinen Namen bekennen und ihm von Herzen dafür danken.

Solchen großen herrlichen Sieg erlanget der HErr auch wunderbarer Weise, stellet sich noch närrischer (wie es die Vernunft ansiehet) dazu, denn David: stirbt des schmähligen Todes am Kreuz, zwischen zween Uebelthätern, als der ärgste Aufrührer und Gotteslästerer; richtet aber dadurch das aus, daß ihm Sünde, Tod und Teufel forthin ausgezogen und erwürget, in Ewigkeit zu Füßen liegen müssen und seinen Christen nicht schaden mögen.

\section{Votum}

Der treue GOtt stehe bei allen, die heute unter dem Kreuz, in Anfechtung wider den Satan streiten, und helfe ihnen und uns ritterlich überwinden, daß wir alle in Christo Jesum, der des Teufels Werk zerstöret hat, selig werden, Amen.

\chapter{Das 18. Kapitel}

\section{Vorrede}

Gott suchet heim mit dem Kreuze, beide, die Gottlosen und die Frommen. Die Gottlosen darum, daß sie ihrer Sünde wegen gestraft werden, andern zum Exempel und Warnung, damit sie sich für dergleichen Sünden hüten, sollen ihnen nicht dergleichen Sündenstrafen folgen. Die Frommen greift Gott an mit dem Kreuze, auf daß sie nicht in Hoffart gerathen und der Gaben wegen, so er ihnen verliehen, sich nicht überheben, und daß er den zukünftigen Sünden also wehre. Dieß lehret St. Paulus in seinem eigenen Exempel und sagt 2 Cor. 12, 7.: Auf daß ich mich nicht der hohen Offenbarungen überhebe, ist mir gegeben ein Pfahl ins Fleisch, nämlich, des Satans Engel, der mich mit Fäusten schlage, auf daß ich mich nicht überhebe. Und St. Petrus 1. Epistel 4. Kapitel V. 12. spricht: Wer am Fleische leidet, der höret auf von Sünden, daß er hinfort, was noch hinterstelliger Zeit im Fleisch ist, nicht der Menschen Lüsten, sondern dem Willen Gottes lebe. Davon höret auch an das 18. Kapitel.

\section{Davids Lob, Lebensgefahr und Heirath}
\quellseite{436}

\section{Summarium}

Ist eine Historia, in welcher wir sehen, wie ein greulich Ding es sei, wo man einmal von der Furcht Gottes und seinem Willen abgetreten ist und also beharret. Denn Saul, da er zum andernmal wider Gott gehandelt hatte, entzog ihm Gott erstlich seinen Geist und ließ ihn plagen mit einem bösen Geist. Derselbe treibet ihn, wie wir hören, daß er David feind wird, und ihn begehrt umzubringen, V. 29. Solche Gefahr aber verhänget Gott über David, daß er nicht stolz noch sicher würde, weil ihn Gott so mit einem herrlichen Sieg begnadet hatte; sondern er weiter Ursache hätte, Gott zu bitten um seine Hülfe und Gnade, und sich seines Worts und Zusagung zu trösten.
\quellseite{437}

\section{Votum}

Gott regiere und behüte uns alle, daß wir durch Sünde den heiligen Geist nicht verjagen, und helfe, daß wir alles Kreuz in Geduld tragen, und uns auch seiner Verheißungen, Beistandes und Hülfe allzeit trösten, Amen.

\chapter{Das 19. Kapitel}

\section{Vorrede}

Der Teufel führet die Menschen aus einer Sünde in die andere. Gott erlöset seine Gläubigen aus einem Kreuz nach dem andern. Dieses soll uns dienen zur Warnung und zu Trost. Zur Warnung, daß, wer einmal dem Teufel gefolget, daß er alsbald Buße thue, auf daß ihn der Teufel nicht in mehr Sünden und denn in ewiges Verderben führe. Zu Trost, daß Gott die Seinen in Gefahr ihres christlichen Berufs erretten und erhalten will. Diesen Trost lehret Gott durch den Propheten Esaias im 54. Kapitel V. 14., da er also saget: Du wirst ferne sein von Gewalt und Unrecht, daß du dich dafür nicht dürfest fürchten; und von Schrecken, denn es soll nicht zu dir nahen. Aller Zeug, der wider dich zubereitet wird, dem soll nicht gelingen. Das ist das Erbe der Knechte des Herrn, und ihre Gerechtigkeit von mir, spricht der Herr. Dies lehret auch das 19. Kapitel.

\section{David begibt sich, das Leben zu erhalten, ins Elend}
\quellseite{438}

a) Bloß nicht daß er nackt gewesen sei, sondern hat die königlichen Kleider abgelegt, und nur gemeine Kleider anbehalten, als ein anderer Mensch. Und fiel nieder den ganzen Tag etc. das ist, er betete mit ihnen, und wenn sie niederfielen, fiel er auch nieder mit ihnen.

\section{Summarium}

Wie ist eine Historia, in welcher man siehet, wie der Teufel nicht feiert, sondern wo er einmal hinter einen Menschen kömmt, treibet er ihn je länger je mehr in größere Sünde. Wie Saul sich hie öffentlich lässet hören, er wolle David auch erwürgen, und nimmt's auch mit der That für, V. 10. 11.

Hie siehet man auch, wie durch wunderbare Mittel Gott den unschuldigen David warnet und errettet, V. 23. Denn Gott will die nicht verlassen, so sich seiner Hülfe in einer guten Sache trösten. Darum muß Saul ehe wieder zum Propheten werden, und seine Knechte auch, ehe dem frommen David etwas von ihnen widerfahren soll.

\section{Votum}

Der treue, barmherzige GOtt lehre und regiere uns, daß uns der Teufel in Sünden nicht verführe; steure auch allen Widerwärtigen seiner Christenheit, und erhalte unter dem Kreuz alle, die seinen heiligen Namen lieben und ehren, Amen.

\chapter{Das 20. Kapitel}

\section{Vorrede}

St. Paulus 2 Tim. 3, 12. sagt: Alle die gottselig leben wollen in Christo Jesu, müssen Verfolgung leiden. Und 1 Cor. 10, 13. spricht er: Gott ist getreu, der euch nicht läßt versuchen über euer Vermögen, sondern macht, daß die Versuchung so ein Ende gewinne, daß ihr könnet ertragen. Danach St. Petrus 1 Epistel 5, 8. sagt: Euer Widersacher, der Teufel, gehet umher, wie ein brüllender Löwe, und suchet, welchen er verschlinge. Und in der Offenbarung Johannis im 12. Kapitel v. 12. spricht eine Stimme im Himmel: Wehe denen, die auf Erden wohnen, und auf dem Meer; denn der Teufel kommt zu euch hinab, und hat einen großen Zorn, und weiß, daß er wenig Zeit hat. Der Spruch St. Pauli ist ein Exempel der verfolgte David. Des Spruchs St. Petri und der Stimme im Himmel ist ein Exempel der Tyrann Saul. Von beiden lautet das 20. Kapitel.

\section{Freundschaft Jonathans und Davids}
\quellseite{439}

\section{Summarium}

Dies ist eine Historia, in welcher man siehet, wie dem frommen David so angst und bange bei den Sachen wird, daß er nirgend seines Leibes und Lebens sicher ist, V. 1. Solches Kreuz und Anfechtung kommet ihm unverdienter Sache; denn ers je um Saul nicht verdienet. Auf daß wir lernen, wie Gott seine Heiligen so mancherlei versucht, und ihnen Ursache genug gibt, daß sie beten und nicht sicher sein sollen. Aber wiederum ists tröstlich, daß ihm Gott so einen treuen Freund bescheret an Jonathan, daß er ihn warnet und für Unglück behütet, V. 4.

a Das ist, du bist nicht ein tapferer, berüchtigter, sondern eines verachteten, geringen und unartigen Weibes Sohn.
\quellseite{440}

\section{Summarium}

Ist ein schrecklich Exempel, in welchem wir sehen, wie der Teufel den König Saul aus einer Sünde in die andere wirft. Bisher hat er David nach dem Leben gestanden, jetzt schießt er nach seinem eignen Sohn, und begehrt ihn auch zu erwürgen. Darum ist es hoch vonnöthen, daß man sich mit allem Ernst in Gottes Furcht bewahre und für Sünden hüte. Denn kommet man einmal darein, so weiß man nicht, wie und wenn man wieder heraus kommet, wie dies Exempel anzeigt.

\section{Votum}

Gott regiere uns, daß wir uns den Teufel aus einer Sünde in die andere nicht führen lassen; und stehe seinen Gläubigen bei in Versuchung und Anfechtung, daß sie alles das in Geduld tragen, und ritterlich anstehen und überwinden, Amen.

\chapter{Das 21. Kapitel}

\section{Vorrede}

Die Liebe und Noth haben kein Gebot. Dieses ist wohlgeredt, man soll es aber auch recht verstehen. Vom Gehorsam gegen menschlichen Satzungen ist es ingemein wahr; aber was Gottes Gesetze anlanget, soll man das verstehen von dem Gehorsam gegen den Statuten und Ordnungen vom jüdischen Gottesdienst, und Ceremonien in der israelitischen Kirche gebräuchlich, und denn von jüdischer Polizei oder gerichtlichen Gesetzen; doch so ferne, da diese mit den natürlichen oder moralischen Gesetzen, das ist, mit den heiligen Zehen Geboten, nicht zustimmen. Denn was in diesen levitischen Kirchen Ceremonien und auch in weltlichen Gerichtordnungen der Juden mit den heiligen Zehen Geboten ingemein einstimmet, das soll auch gehalten und gethan werden. Dieser Lehre ist ein Exempel das 21. Kapitel.

\section{David hungert, fliehet und verstellet sich}

a Zeug in der Schrift heißt ein jeglicher Leib ein Zeug, wie auch Act. 9, 15. Christus von St. Paulo sagt: Er ist
\quellseite{441}

\section{Summarium}

Ist eine Historia, wie David in der Noth, da er sonst nichts hatte zu essen, die heiligen Brode isset, die doch sonst niemand denn die Priester essen durften, V. 6. Solches deutet Christus im Evangelio dahin, daß man das Gesetz so hart nicht spannen, sondern allweg auf die Noth leben und um der Noth willen etwas nachgeben soll. Auf daß wir nicht thun, wie die Phariseer, die hätten eher einen Menschen verderben lassen, ehe sie am Sabbath ihm geholfen hätten. Und die Papisten ließen eher einen Menschen sterben, ehe sie ihm am Freitag Fleisch geben. So doch Gott solches nirgends verboten hat.

\section{Votum}

Der treue Gott nehme sich unserer Noth an, als ein treuer, lieber Vater, und regiere uns, daß wir unsern Nächsten in seiner Noth auch nicht verlassen, sondern ihm allezeit mit Rath und That gern und treulich dienen. Amen.

\chapter{Das 22. Capitel}

\section{Vorrede}

Der Teufel feiret nicht, sondern wo er einmal hinter einen Menschen kömmt, so treibet er ihn je länger je mehr in andere und größere Sünde. Dergleichen ist es hoch vonnöthen, daß man sich mit allem Ernst und Fleiß in Gottes Furcht bewahre, für Sünden hüte, und Gott um Hülfe anrufe, daß der Teufel unser nicht mächtig werde, und uns aus einer Sünde in die andere, zu unserm Verderben, nicht führe. Dieses lehret der HErr Christus, da er vermahnet Matth. 26, 41.: Wachet und betet, daß ihr nicht in Anfechtung fallet. Dies lehren auch beide heiligen Apostel des HErrn Christi, Petrus und Paulus. Petrus saget 1. Epistel 5. Kapitel V. 8.: Seid nüchtern und wachet, denn euer Widersacher, der Teufel, gehet umher wie ein brüllender Löwe, und suchet, welchen er verschlinge; dem widerstehet fest im Glauben. Und St. Paulus Ephef. 6, 10. spricht: Seid stark in dem HErrn und in der Macht seiner Stärke; ziehet an den Harnisch Gottes, daß ihr bestehen könnet gegen die listigen Anläufe des Teufels. Und betet stets in allem Anliegen mit Bitten und Flehen im Geist. Dies lehret auch das 22. Kapitel.

\section[Saul läßt durch den Verräther Doeg fünf und]{Saul läßt durch den Verräther Doeg fünf und achtzig Priester tödten}
\quellseite{442}

\section{Summarium}

Ist eine schreckliche Historia, wie der böse Geist den König Saul immer in mehr und greulichere Sünde treibet, daß er ohn alle Ursache 85 Priester und eine ganze Stadt mit Weib und Kind, mit allem Vieh, erwürgen lässet, V. 18.; und solchen Mord noch mit Lügen schmücket: David stelle ihm nach Leib und Leben, V. 13. So doch David für ihn, ohn alle Ursache, weder Leibes, noch Lebens sicher war. Also siehet man, wie der Teufel ein Lügner und Mörder ist. Darum soll man sich mit allem Fleiß für
\quellseite{443}
Sünden hüten, und stets bitten, daß uns Gott nicht in Versuchung führen lasse.

\section{Votum}

Gott, der Herr Zebaoth, kenne des Teufels List und Gewalt, und regiere und stärke uns in unserer großen Schwachheit, daß wir wider des Teufels Anfechtung streiten, und uns von ihm ja nicht verführen lassen, Amen.

\chapter{Das 23. Kapitel}

\section{Vorrede}

Eines jeden Beruf, Handel und Wandel, dienet Gott und dem Nächsten, so man sie in Gottes Furcht thut, und nach Gottes Befehl und Willen; und je größere Noth darinnen fürfället, je näher ist der treue Gott mit Rath und Hülfe. Daher gehöret, daß Gott selber sagt im 91. Psalm V. 15.: Ich bin bei ihm in der Noth, ich will ihn heraus reißen und zu Ehren machen. Denn Gott hat auch mancherlei Mittel, der Gottlosen Anschläge wider die Gottseligen zu verhindern. Also gehöret auch hieher, daß der heilige Apostel Paulus sagt, 2 Cor. 4, 8. 9.: Wir haben allenthalben Trübsal, aber wir ängsten uns nicht; uns ist bange, aber wir verzagen nicht; wir leiden Verfolgung, aber wir werden nicht verlassen; wir werden unterdrückt, aber wir kommen nicht um. Davon höret auch an das 23. Kapitel.

\section[David wird von neuem verrathen, verfolgt, und wunderbarlich]{David wird von neuem verrathen, verfolgt, und wunderbarlich errettet}

\section{Summarium}

Ist eine schöne Historia, wie David, ehe er
\quellseite{444}
wider die Philister, seinem Volk zu gut, ausgezogen, zuvor den HErrn fraget, ob ers thun soll, V. 2. Denn ohne Beruf soll niemand nichts fürnehmen. Darnach ist auch ein fein Exempel, wie Gott den frommen David warnet für den untreuen und undankbaren Bürgern zu Kegila, V. 11. Auf daß wir lernen, obgleich Gott den Seinen mancherlei Kreuz und Anfechtung zuschicket, so schaffet er doch alleweg Hülfe und Rath daneben.

a Sela Mahelkoth) das heißet Scheidefels.

\section{Summarium}

Ist eine Historia, wie der fromme David von den Siphitern ausgespähet und dem König Saul verrathen ward. Aber Gott schaffts wunderbarlich, da ihn Saul umringet, und am gewissen hatte, fallen ihm die Philister ins Land, daß er muß abziehen, V. 27. Gleich wie zu unserer Zeit die Feinde des Wortes Gottes oft Arges im Sinn gehabt haben; aber Gott hats bisher allezeit also geschicket, daß sie von ihrem Fürnehmen haben ablassen müssen. Er wirds auch forthin thun, wenn wir mit Beten getrost anhalten.

\section{Votum}

Der barmherzige Gott, der treue Menschenhüter, rathe und helfe heute auch allen frommen Christen, die in Noth sind, und steure allem bösen Fürnehmen der Feinde seiner Kirchen, damit sie erhalten werden wider alle Pforten der Hölle, Amen.

\chapter{Das 24. Kapitel}

\section{Vorrede}

Rachgier ist eine Sünde, die Gott erzürnet, das Gewissen beschweret und ihre eigene Sache verderbet; dagegen, wer verhören und verzeihen kann, der gefällt Gott und den Menschen wohl, da wird Ehre aus, bleibet ungebissen und verhütet viel Schaden. Derwegen wer ein rechter Christ sein will, soll sich selbst überwinden, Geduld haben und sanften Muth, und auch dem Feinde thun das Gut. Dies lehret der HErr Christus Matth. 5, 44., da er also spricht: Ich sage euch, liebet eure Feinde, segnet, die euch fluchen, thut wohl denen, die euch hassen, bittet für die, so euch beleidigen und verfolgen, auf daß ihr Kinder seyd eures Vaters im Himmel. Dies unterrichtet auch St. Petrus 1. Epistel am 3. Kapitel V. 9., da er sagt: Vergeltet nicht Böses mit Bösem. Und St. Paulus Röm. 12, 19. vermahnet: Rächet euch selbst nicht, meine Liebsten, sondern gebet Raum dem Zorn (Gottes). Denn es stehet geschrieben, die Rache ist mein, ich will vergelten, spricht der HErr. Davon höret auch an das 24. Kapitel.

\section{David kommt bey Saul in der Höhle}
\quellseite{445}

a Sein Fuß zu decken) so pflegt es die heilige Schrift, daß sie Füße decken heißt, auf das heimliche Gemach gehen.

\section{Summarium}

Es ist eine sehr schöne Historia, wie David, da er seinen Feind, den König Saul, schon in seinen Händen hatte, dennoch sein schonet und sagt, er wolle an den Gesalbten des HErrn die Hand nicht legen, V. 7. Denn ob ihn Gott gleich zum König lassen salben, hat er ihm doch nicht befohlen, daß er Saul tödten und das Königreich ihm selbst nehmen sollte, sondern warten, bis ihm Gott in die Hände gebe. Wiederum thut David recht, daß er Saul seine Unschuld anzeiget, und ihn um sein unbillig Fürnehmen strafet, V. 11. Denn großen Herren soll man ihre Sünde eben so wenig billigen und verschweigen, als armen und geringen Leuten.

\section{Votum}

Verleihe GOtt, daß wir alle, Obrigkeit und Unterthanen, die Sündenstrafen aus GOttes Wort annehmen, und daß wir unsere Feinde auch mit christlicher Geduld überwinden, und ihnen feurige Kohlen auf ihr Haupt sammeln, Amen.
\quellseite{446}

\chapter{Das 25. Kapitel}

\section{Vorrede}

Beide die Gottlosen und Frommen haben zu Sünden Zuneigung, aber mit großem Unterschied. Die Gottlosen folgen ihren bösen Zuneigungen und Lüsten, dem Teufel zu gefallen und ihnen selbst zu Schaden. Die Frommen aber widerstreben den bösen Zuneigungen und Lüsten, Gott zu Ehren, und daß sie nicht, weder ihnen selbst, noch andern, Unglück zurichten. So lasset nun die Sünde nicht herrschen in eurem sterblichen Leibe, ihr Gehorsam zu leisten in ihren Lüsten. Auch begebet nicht der Sünden eure Glieder (vernehmet nicht allein die Augen, Ohren, Hände und Füße, sondern auch die innerlichen Sinne, das Herz und seine Gedanken) zu Waffen der Ungerechtigkeit (verstehet, Unrechtes zu thun); sondern begebet euch selbst Gott und eure Glieder (vernehmet aber, nicht allein die äußerlichen, leiblichen Gliedmaßen, sondern auch alle innerlichen Sinne und Kräfte) Gott zu Waffen der Gerechtigkeit (vernehmet, daß ihr im Glauben Gott neuen Gehorsam leistet, in aller Tugend und Ehrbarkeit): und diese Vermahnung thut der heil. Geist an alle Christen, durch den heiligen Apostel Paulus an die Römer, am 6. Kapitel V. 12. 14. Dies lehret auch der erste Theil des 25. Kapitels.

\section[Samuels Tod. Nabals Thorheit. Vorsichtigkeit Abigails. Davids Werbung]{Samuels Tod. Nabals Thorheit. Vorsichtigkeit Abigails. Davids Werbung um sie}

\section{Summarium}

Dies ist erstlich ein Exempel an Nabal, wie bisweilen grobe Leute ihnen selbst alles Unglück mit ihrer Grobheit zurichten, V. 10.

Darnach ist ein Exempel, wie David sich den Zorn übergehen läßt, und gedenkt solche Grobheit und Unfreundlichkeit an Nabal zu rächen, und sein eigener Richter zu sein, V. 13. Aber solches war Sünde und Unrecht. Denn obwohl Nabal hatte unrecht gehandelt, sollte doch David nicht sein selbst Richter sein. Also begegnet es auch den heiligen Leuten, daß sie zu jach und geschwind sind, und Unrecht zu thun fürnehmen. Aber sie lassen sichs bald gereuen, und hören auf, da die andern fortfahren, und ihnen nichts abbrechen.
\quellseite{447}

\section{Summarium}

Ist ein schön Exempel, wie Gott Mittel schicket, daß David nicht fortfähret und sich an Nabal rächt, wie er im Sinn hatte. Und hie siehet man, wie man mit Freundlichkeit und guten Worten oft eine Sache führen und Unglück verhüten kann. Darum soll jedermann sich befleißigen, daß er sich nicht bald lasse bewegen zu Zorn und Rache, daß ihn darnach gereuen möchte. Sonderlich aber ists fein, daß Abigail zu David sagt, wo er ihm selbst geholfen und sich an seinem Feind gerochen hätte, würde ihm solche Sünde ein Anstoß und Aergerniß gewesen
\quellseite{448}
sein, wenn er zum Regiment wäre kommen, V. 13. Denn Gott läßt solche Sünde nimmer ungerochen abgehen.

\section{Votum}

GOtt vergebe es allen, die ihre Hand von gottfürchtigen Herzen abziehen, und welche die verlassen, denen sie helfen sollen; und regiere uns, daß wir unserm Fleisch und Blut zu Sünden nicht folgen, sondern GOtt und seinem Wort folgen, in allem Guten, unserm Nächsten zu Nutz, Amen.

\section{Der andere Theil des 25. Kapitels}

\section{Vorrede}

Unsere lieben alten Deutschen haben recht gesagt: Wer jach, dem gebrach; wer vertrug, der hat genug; wer schweiget, der bleibet; Gott richt't, wenn niemand spricht. Denn Gott sagt selber 5 Mos. 32, 35.: Mein ist die Rache, ich will vergelten. Derwegen für Nachgier sich männiglich hüten soll. Auch sind dieses der alten Deutschen gute Sprüchwörter gewesen: Der hat genug, der sich begnügen läßt. Laß dir begnügen, so ist dir Gott und die Welt hold. Aus kleinen Brünnlein trinket man sich auch satt, sowohl als aus großen. Laß dir nicht so viel gelieben, so wird dichs nicht betrüben. Denn also sagt der Apostel Paulus 1 Tim. 6.: Es ist ein großer Gewinn, wer gottselig ist und lässet ihm genügen. Derwegen für Geiz sich auch jedermann soll hüten. Dies soll ferner erinnern der andere und letzte Theil des 25. Kapitels.

\section{Summarium}

Ist ein fein Exempel, wo man Gott läßt Richter sein und seiner Hülfe erwartet, daß Gott nicht außen bleibt, sondern die Untugend straft, wie man verdienet hat. Wie man siehet, daß Gott den David an Nabal rächet, und ihm, wie er verdienet hatte, vergilt, daß er sich zu Tode kümmeret und stirbt, V. 38. Solches Exempel gehet aber sonderlich auf die Reichen und Geizigen, so um Geizes willen andern fürhalten oder kärglich bezahlen, was sie ihnen schuldig sind, und man ihnen abverdienet hat. Darum hat dieser Nabal über die Maßen viele Brüder in allen Städten und Dörfern. Und wir sehen, wie es letztlich mit ihnen hinaus gehet, daß sie dahin fallen und sterben, und selbst ihres Guts nicht genießen noch froh werden, welches sie andern nicht gönnen; bisweilen aber wird ihnen bei alle ihrem Gut so angst und bange, daß sie sich selbst hängen oder sonst umbringen. So ein schädliches, greuliches Laster ist der Geiz.

\section{Votum}

Der treue GOtt und Vater unsers HErrn JEsu Christi segne ferner alle, die von seinem Segen andern mittheilen, und regiere uns, daß wir uns für sündlichem Zorn, Nachgier, Geiz und allen andern Sünden hüten, und aller Sünden Strafe entfliehen, Amen.

\chapter{Das 26. Kapitel}

\section{Vorrede}

Gleich wie fromme Herzen ihre Feinde mit Geduld und Wohlthaten überwinden, daß sie sich schämen müssen und ihre Uebelthat bekennen; denn das ist allein die rechte Form und Weise, an Feinden sich zu rächen: also stellen sich die Heuchler freundlich, dieweil es ihre Noth erfordert, wenn die aber ein wenig vorüber ist, da hat ihre Freundlichkeit auch ein Ende, und wird also offenbar, daß es eine gefärbte Liebe und Versöhnung sey, dahinter nichts denn Galle, Blut, Mord und Verderben steckt. Daher sagt

a Neminem invitat, nec pauperes curavit. Das ist, er hat Niemand eingeladen, und sich nicht um die Armen bekümmert.
\quellseite{449}
Salomon in Sprüchen am 29. Kapitel: Die Blutgierigen hassen den Frommen, aber die Gerechten suchen seine Seele. Derwegen niemand sich selbst an seinen Feinden soll rächen, sondern soll durch Gott und die weltliche Obrigkeit sich rächen lassen, und soll sich auch das Böse nicht lassen überwinden, sondern das Böse soll er mit Gutem überwinden, Röm. 12, 21., das ist, durch Geduld, Sanftmuth und Wohlthat gegen seinen Feinden. Dies soll auch erinnern das 26. Kapitel.

\section{David nimmt Sauls Spieß und Wasserbecher}
\quellseite{450}

\section{Summarium}

Dies ist eine Historia, wie David zum andern mal seinen Feind in seinen Händen hat, und dennoch sich enthielt, sein schonet und nicht sein selbst Richter wird, V. 10. Solches läßt Gott dem Saul zu gut geschehen, daß er den unbilligen Haß und Zorn gegen David fallen lassen, und im Werk lernen und erkennen soll, daß ihm David kein Leid zu thun begehrt. Aber die Gottlosen verachten entweder solche Warnung gar, oder haben nur eine kleine Galgen-Reue, thun hernach wie vor, wie man hie an Saul siehet.

\section{Votum}

Gott steure und wehre allen falschen Herzen, und regiere uns, daß wir treue, ungefärbte, beständige Liebe gegen unsern Freunden und Feinden tragen, daß Gott zu Ehren und ihnen zur Besserung, Nutz und Wohlfahrt geistlich und leiblich diene. Amen.

\chapter{Das 27. Capitel}

\section{Vorrede}

Der heilige Apostel Petrus 1. Epistel 3. Capitel V. 9. vermahnet und spricht: Vergeltet nicht Böses mit Bösem, sondern dagegen segnet, das ist, seid nicht rachgierig, sondern sanftmüthig, und um gemeines und eigenes Friedens willen, vertraget das Unrecht. Und setzet eine gewaltige Ursache dazu, warum solches geschehen soll, nämlich: Und wisset, daß ihr dazu berufen seid, daß ihr den Segen beerbet; das ist, weil ihr von Gott gesegnet seyd in dem gesegneten Samen Abraham, Jesu Christo, daß ihr den ewigen Segen mit seiner Herrlichkeit und Freuden habet und ererbet, ist es billig und recht, daß ihr wiederum segnet, Gutes wünschet und betet, denen, die euch fluchen. Er vermahnet auch dazu mit angehängtem herrlichen Trost, da er saget: Wer ist, der euch schaden könnte, so ihr dem Guten nachkommet? Und ob ihr auch leidet um Gerechtigkeit willen, so seyd ihr doch selig. Als wollte er sagen: Denn eure Widersacher mit ihrer Bosheit und ungerechter Gewalt euch Gottes Gnade, Hülfe und Segen, und eure Frömmigkeit nicht nehmen können. Davon höret auch das 27. Kapitel.

\section{Davids Wohnung in Ziklag}
\quellseite{451}

\section{Summarium}

Ist eine Historia, wie David sich dem Saul nicht zu weit vertraut, sondern aus dem Wege gehet, V. 1. Und nichts desto weniger seinem Volk Israel zu gut die umliegenden Heiden ausrottet und verderbet, welche Gott dem Volk Israel auszurotten geboten hatte, V. 8. Auf daß man lerne sich überwinden, und nichts Böses mit Bösem bezahle. Denn weil Saul dem David nachstellet, und ihm das Volk zu solchem half, hätte David wohl Ursache für der Welt gehabt, sich an ihnen, als seinen Feinden, zu rächen. Aber er thuts nicht, sondern wartet, bis es Gott räche, und thut ihnen dieweil alles Gutes.

\section{Votum}

Der allwissende und gütige Gott erleuchte und regiere uns, daß wir uns in dieser untreuen Welt auch wohl vorsehen, und helfe, daß wir als seine wohlgerathenen und gehorsamen Kinder seinem Exempel nachfolgen, und Freunden und Feinden gern und willig dienen, Amen.

\chapter{Das 28. Capitel}

\section{Vorrede}

Ein jedes frommes Herz und gläubiger Mensch soll täglich gedenken an König Saul, also, daß es sich dabei erinnere des großen Zorns Gottes wider die Sünde, und der erschrecklichen Strafen, die aus Gottes gerechtem Gericht allen Unbußfertigen folgen; und darnach, daß ein jedes fleißig bete, Gott wolle ihn nicht verwerfen unter die Gefäße des Zorns und unter die Landschaden des menschlichen Geschlechts, sondern daß er ihn machen wolle zum Werkzeuge der Gnaden und Barmherzigkeit, ihm, Gott dem Herrn, zu Ehren und dem Nächsten zu Heil und Wohlfahrt. Und vermögen soll aller Christen tägliches Gebetlein sein, daß sie viel und oft mit dem königlichen Propheten David seufzen sollen aus dem 119. Psalm V. 133.: Herr, laß meinen Gang gewiß sein in deinem Wort, und laß kein Unrecht über mich herrschen. Dies lehret auch das 28. Kapitel.

\section[Saul suchet Rath bei einem Zauberweibe, und wird]{Saul suchet Rath bei einem Zauberweibe, und wird vom Gespenst erschreckt}

ᵃ Das erzählet die Schrift darum, auf daß sie warne jedermann, daß er das nachfolgende Beispiel von Samuel recht verstehe, und wisse, daß Samuel todt sei, und solches der böse Geist mit der Zauberin und Saul redet und thut, in Samuels Person und Namen.

ᵇ Das Licht ist, das auf dem Brustlaz des Priesters war.
\quellseite{452}

c Götter) das ist, Richter. 2 Mos. 22, 8. Und der seidene Rock ist der priesterliche Rock.

\section{Summarium}

Ist erstlich ein Exempel, wie Saul immer in größere Sünde fällt. Er ist bisher Gott mehr denn einmal ungehorsam gewesen, hat den unschuldigen David mehr denn einmal umbringen wollen, nach seinem eigenen Sohn hat er einen Spieß geschossen. Sie sehen wir, wie er von Gott abfället und des Teufels Hülfe und Rath bei einer Wahrsagerin suchet, V. 7. Sogar ist kein Aufhören von Sünden, wo man einmal anhebet und nicht Buße thut.

Sonderlich aber ist hie zu merken das Exempel Sauls, welcher, wie hie der Text anzeiget (ohne Zweifel aus Rath und Befehl Samuelis) alle Wahrsager und Zeichendeuter aus dem Lande vertrieben hatte, V. 3. Denn weltlicher Obrigkeit gehörets ihres Amts halben, daß sie solche und andere dergleichen öffentliche Abgötterei und Aberglauben nicht leiden, sondern mit Ernst strafen und abschaffen solle.

\section{Summarium}

Es ist erstlich ein Exempel, in welchem man siehet, wie der Teufel bisweilen kann zukünftige Dinge sagen. Aber nichts desto weniger soll man solcher Weissagung nicht brauchen. Denn es ist eine Sünde wider das erste Gebot, da uns geboten wird, wir sollen uns nur eines Gottes halten, mit Glauben Hülfe und Rath suchen, und allen andern.

Daß aber hie die Schrift sagt, Samuel sei auferwecket und habe mit Saul geredet, V. 15., ist zu verstehen, nicht daß Samuel gewesen sei, sondern daß Saul solches also gemeinet und verstanden habe. Denn so mächtig ist der Teufel selbst nicht, ich geschweige denn ein Mensch, der des Teufels Hülfe braucht, daß er einen verstorbenen Menschen aus seiner Ruhe erwecken und aus dem Schooß Abraham, wie es die Schrift heißet, rücken könnte. Darum ist es gewiß, daß solches nur ein Teufelsgespenst gewesen ist, und
\quellseite{453}
nicht Samuel, der in seiner Ruhe und Schooß Abrahams blieben ist. Die Schrift aber heißt ein solch Gespenst den Samuel darum, daß Saul und die Wahrsagerin es für den Propheten Samuel gehalten haben.

\section{Votum}

Der treue Gott erbarme sich unser, behüte uns für Abgötterei, Zauberei und andern Sünden, und verleihe, daß wir uns allezeit nach seinem Wort und Willen, gegen ihm und gegen allen Menschen christlich erzeigen, bis zu unserm seligen Todesstündlein, Amen.

\chapter{Das 29. Kapitel}

\section{Vorrede}

Wie einer, den die Pestilenz angestoßen, da er ihm in 24 Stunden durch ordentliche Mittel nicht Rath schaffen läßt, sich selbst verwahrloset, daß solch Gift durch alle Adern, Geblüt und ins Herz, hinein dringet, und daß ihm darnach nicht kann geholfen werden, sondern er muß sterben: also müssen auch die, welche der geistlichen Pestilenz, der Sünden, in Zeit der Gnaden und am Tage des Heils nicht zuvor kommen, ewig daran sterben und verderben. Dies lehret St. Paulus an die Galater im 5. Kapitel V. 19. 21. und saget: Offenbar sind die Werke des Fleisches, als da sind Ehebruch, Hurerei, Unreinigkeit, Unzucht, Abgötterei, Zauberei, Feindschaft, Hader, Neid, Zorn, Zank, Zwietracht, Rotten, Haß, Mord, Saufen, Fressen und dergleichen, von welchen ich euch zuvor gesagt und sage noch zuvor, daß die solches thun, werden das Reich Gottes nicht erben. Derwegen laßt uns alle in täglicher Buße leben; und dazu soll uns auch bewegen das Gnaden- und Straf-Exempel. Davon das 29. Kapitel.

David wird von den Philistern zurückgeschickt.

\section{Summarium}

Ist eine feine Historia, weil David mit dem Könige der Philister mußte ziehen, wie Gott Mittel schaffet, daß er mit gutem Willen des Königs wieder abzeucht, V. 4. Nun ist aber kein Zweifel, wo David bei der Schlacht wäre blieben, er würde wider die Philister, mit denen er ausgezogen war, seinem Volk Israel zu gut gestritten haben, wie sie sichs denn auch besorgen, und ihn darum wieder abschaffen. Aber Gott wollte solches Fürnehmen Davids dem Saul und ganzen Israel zu gut nicht lassen kommen, sondern wollte sie um ihre Sünde, wie sie verdienet hatten, strafen. Darum verschafft er solches Mittel, daß David daheim bleibt, und die Philister das Volk Israel schlagen.
\quellseite{454}

\section{Votum}

Der barmherzige GOtt strafe uns ja nicht, wie wir leider auch verdienet haben, mit unsern Sünden; und behüte uns, daß wir uns an niemand versündigen, daß also unser Leben ihm zu Lob und unserm Nächsten zu Nutz diene, Amen.

\chapter{Das 30. Kapitel}

\section{Vorrede}

Die Gottseligen sind selten ohne Kreuz, ja ihr Kreuz wächst und nimmt täglich zu, auf daß auch ihr Glaube, Geduld, Hoffnung und Standhaftigkeit wachse und zunehme; ja zugleich Gottes Liebe gegen sie; denn welchen Menschen der HErr lieb hat, den züchtiget er, Ebr. 12, 6. Je größer aber nun die Züchtigung ist, je größer ist auch die Liebe Gottes; wie der nützliche und fröhliche Ausgang des Kreuzes aller Gläubigen beweiset. Denn die Errettung aus dem Kreuz viel herrlicher ist, denn alles vorgehende Leiden der Gottseligen, auf daß sie also auch ihrem Nächsten denn aus christlicher Liebe und brüderlicher Treue rathen, dienen und helfen können. Denn solches der Liebe Art ist, als St. Paulus lehret 1 Cor. 13, 4. 7., da er sagt: Die Liebe ist freundlich, sie suchet nicht das Ihre, sie lässet sich nicht erbittern, sie trachtet nicht nach Schaden, sie verträget alles, sie duldet alles. Dessen ist ein schön Exempel David, von dem das 30. Kapitel.

\section[David sieget wider die Amalekiter, und theilet den]{David sieget wider die Amalekiter, und theilet den Raub aus}
\quellseite{455}

\section{Summarium}

Ist eine Historia, wie dem David, weil er mit den Philistern wider Israel zeucht, die Feinde einfallen, Weib und Kind sammt aller Habe entführen und die Stadt Ziklag verbrennen. Aus solchem kömmt David zu einem andern Unrath, daß ihn sein eigen Volk will steinigen, V. 6. Aber er hält sich zu seinem Gott, und tröstet sich desselbigen, und jaget aus Gottes Befehl den Feinden nach und erobert alles wieder, V. 10. Auf daß wir lernen, wie Gott mancherlei Anfechtung den Seinen zuschickt, aber auch getreulich wiederum heraus hilft, wenn man sich seiner Gnade und Hülfe tröstet.

\section{Summarium}

Hie ist ein fein Exempel, wie David seinen Leuten zur Versöhnung räth, und Zwietracht durch billige Mittel verhütet, daß er den Raub gleich austheilt. Denn niemand soll Geld oder Gut besser und höher achten, denn Friede und Einigkeit. Wer aber solches nicht thut, und um Geld und Gut willen sich in Unfriede begibt, der kommt endlich auch ums Geld, wie die Exempel täglich für Augen sind.

\section{Votum}

Unser lieber GOtt und treuer Vater helfe heute auch, zu Trost und Rettung allen christlichen Kreuzträgern; und regiere uns, daß wir jedermann in Lieb und Leid, und nach unserm höchsten Vermögen, und auch zu Friede und Einigkeit dienen, Amen.

\chapter{Das 31. Kapitel}

\section{Vorrede}

Gott läßt wahrhaftig die Sünden nicht ungestraft, ob er gleich, als ein sanftmüthiger, barmherziger, gnädiger und langmüthiger Gott, mit der Strafe eine Zeit lang aufseucht und verschonet. Es strafet aber Gott, der Herr, nicht allein Vater und Mutter, die seine Gebote übertreten und sündigen, sondern er suchet deshalb auch ihre Kinder mit Ungnaden heim, daß sie also der Eltern Mißhandlung entgelten müssen. Dieses lehret Gott 2 Mose im 20. Kapitel V. 5., da er sagt: Ich, der Herr dein Gott, bin ein eifriger Gott, der da heimsuchet der Väter Missethat an den Kindern bis in das dritte und vierte Glied, die mich hassen. Demegegen die Eltern gottesfürchtig sein und ehrbar leben sollen, nicht allein, daß Gott sie nicht ihrer Sünden wegen strafen dürfe, sondern auch darum, daß sie ihren Kindern seinen Segen und nicht seinen Fluch und Strafe aufgeben und lassen. Dieses lehret auch das 31. und letzte Kapitel.

\section[Sauls letzter Krieg wider die Philister, darinnen er]{Sauls letzter Krieg wider die Philister, darinnen er sich selbst erstochen}
\quellseite{456}

\section{Summarium}

Ist eine Historia, in welcher man siehet, wie es endlich mit dem Saul ein Ende nimmt, daß er mit seinen drei Söhnen erstochen und das Volk Israel dazu geschlagen wird, V. 4. Darum sollen solche Exempel uns dahin halten, daß wir Gott fürchten und Tag und Nacht mit Fleiß beten, daß er uns für Sünden bewahren wolle, auf daß wir nicht endlich in dergleichen Strafe auch fallen, und so jämmerlich sammt den Unsern umkommen. Denn es ist nicht ein Geringes, daß solcher Jammer dem Saul, der so große und hohe Gaben gehabt, widerfahren, und er so greulich und tief gefallen, und die Strafe über alle seine Nachkommen gangen ist. Der allmächtige, gnädige Gott wolle mit seinem heil. Geist uns leiten und führen, daß dergleichen uns nicht auch begegne.

\part{Das zweite Buch Samuelis}

\chapter{Das 1. Kapitel}

\section{Vorrede}

Gott, der HErr, gönnet der Obrigkeit ein Stück seines Titels und Namens, und nennet sie Götter, als im Ps. 82, 1., und darum, daß sie göttlich Amt führen, und daß er sie will geehret haben. Derwegen die Unterthanen ihrer ordentlichen Obrigkeit nicht allein gebührliche, äußerliche Ehrerbietung mit Gebehrden und Worten erzeigen, sondern sie auch von Herzen in Ehren halten sollen; und wenn die Personen, welche in ordentlichem Amt und Regiment sitzen und den Stand der Obrigkeit führen, als Menschen auch ihre Gebrechen haben und nicht alle Dinge zurecht machen, sollen fromme Unterthanen wegen derselben ihrer Gebrechen Geduld und Mitleiden tragen, und sollen für die Obrigkeit beten, so oft sie zusammen kommen, insgemein und darnach auch insonderheit daheim, und was sie Gott zu Ehren und dem Regiment zum Besten thun, mit Dankbarkeit rühmen und preisen. Davon höret auch an das 1. Kapitel des andern Buchs Samuelis.

\section{Davids Traurigkeit über den Tod Sauls und Jonathans}
\quellseite{457}

\section{Summarium}

Ist erstlich eine Historia, wie David den Amalekiter erwürgen lässet, der ihm ansaget, wie er Saul erstochen hätte, V. 15. Auf daß wir lernen, wie der fromme David so hoch geachtet habe, daß Gott den Saul zum Könige gemacht hatte, und wir demselbigen Exempel nach die Obrigkeit auch lernen hoch halten und uns hüten, daß wir nicht die Hand an sie legen. Denn, so sie unrecht thut, soll man Gottes Urtheil und Strafe erwarten, und nicht wider Gottes Ordnung unser selbst Richter sein und unsern Schaden selbst rächen.

a Bogen) so heißt dies Lied, wie auch bei uns etliche Lieder Namen haben.
\quellseite{458}

\section{Summarium}

Ist ein fein Exempel, wie David den König Saul sammt seinen Söhnen beweinet, und befiehlet, daß man ihrer durch ein öffentlich Lied gedenken soll, als tapferer Helden, welche dem ganzen Lande zu gut ihr Leib und Leben wider die Feinde gewaget, und das Land gern geschützt und Frieden erhalten hätten. Denn dies sind die rechten fürstlichen Werke, welche man an großen Herren preisen und loben, und um derselbigen willen sie und die Ihren ehren und hoch halten soll. Wie es David hie auch nennet und saget, es sei der Schmuck, die schönen Kleider, die güldnen Kleinode, welche der König Saul den Töchtern Israel, das ist, den Städten und dem ganzen Lande habe angehängt, und Leib und Leben endlich darob gelassen. Solche Klage ist dem David ein rechter Ernst, denn er siehet, was für ein Ende folget, wo die Regenten dahin fallen.

\section{Votum}

Der barmherzige GOtt regiere, segne und erhalte unsere liebe Obrigkeit, daß sie christlich und wohl ferner regiere; und verleihe, daß wir sie auch herzlich ehren, ihm zum Gehorsam und uns zu gutem Gewissen, Amen.

\chapter{Das 2. Kapitel}

\section{Vorrede}

Vier Sündenstrafen und Plagen hat Gott in der Schrift namhaftig gemacht. Die erste ist die Pestilenz, die andere ist theure Zeit, die dritte ist Krieg, die vierte ist allerlei böse Wesen, als Löwen, Wölfe, Schlangen und dergleichen. Unter diesen ist Krieg und Unfriede die größte Plage und Strafe Gottes, als die da zerstöret und wegnimmt die Religion, weltlich und häuslich Regiment, daß Theurung und Pestilenz, die auch grausame und schädliche Krebse und Wesen Gottes sein, dennoch nicht zu vergleichen sind mit dem Kriege und was im Kriege vorläuft. Derwegen wer danken kann, der danke Gott täglich und herzlich, daß er unseres Vaterlandes so lange Zeit mit Krieg und andern wohlverdienten Strafen gnädiglich verschonet hat; und laßt uns bitten, daß er uns ferner den lieben Land- und Stadtfrieden erhalten, und nöthige Wohlfahrt zu Leib und Seele geben wolle, Amen. Ferner höret an das 2. Kapitel.

\section{Salbung Davids zum Könige. Jeboseths Aufruhr}
\quellseite{459}

a Helkath-Hazurim) das heißt, das Acker der Besten oder Helden.

\section{Summarium}

Dies ist erstlich ein Exempel, wie David nichts fürnimmt, er habe denn zuvor, nachdem es Gott befohlen hat, bei Gott Rath gesucht und Befehl empfangen, V. 1. Darum ordnet es Gott weiter mit ihm, daß ihn das Haus Juda, nämlich der Stamm, darin er geboren war, zum Könige annimmt und salbet, V. 14. Aber es stößet sich hie noch einmal; denn die andern elf Stämme werfen Isboseth, des Sauls Sohn, zum Könige auf, daraus denn Mord und Krieg sich erhebt, wie allweg geschiehet, wo Uneinigkeit in Regimenten fürfället.

\section{Summarium}

Ist ein fein Exempel, wie Abner, nachdem sich Uneinigkeit zwischen der zweien Königen Unterthanen erhoben hatte, zum Frieden räth und Joab ihm solches gefallen lässet und willig abwendt, V. 28., obgleich sein Bruder Asahel in dem Lärmen von Abner erstochen ward, V. 23. Denn in solchen Sachen soll man allewege mehr auf den gemeinen, denn eigenen Schaden sehen, und ums gemeinen Friedens willen des eigenen Schadens gern vergessen.

\section{Votum}

Unser lieber Herr Gott verleihe uns Friede gnädiglich zu unsern Zeiten, und regiere uns, daß wir ein geruhiglich und stilles Leben führen in aller Gottseligkeit und Ehrbarkeit, und in unserm Beruf all unser Thun und Lassen nach seinem Wort anstellen, Amen.
\quellseite{460}

\chapter{Das 3. Kapitel}

\section{Vorrede}

Gleich wie Gott den Verzug der Strafe pfleget mit Schwerheit der Strafe einzubringen, daß ist, je länger Gott verzeucht, die Gottlosen zu strafen, je greulicher er sie darnach, wenn er seine Zeit ersiehet, zu strafen pflegt: also pflegt er auch den Verzug seiner Hülfe einzubringen, und je länger er uns zu helfen verzogen hat, je reichlicher er uns, wenn er seine Zeit ersiehet, zu helfen pfleget. Derwegen sollen wir den Verzug der Hülfe Gottes nicht also verstehen, als wollte Gott uns ganz und gar nicht helfen, und sollen in keinem Wege deshalb das Gebet, Glauben und Hoffnung fahren lassen, sondern sollen vielmehr also thun, wie uns der Prophet Habakuk vermahnet Kapitel 2. V. 3.: Wenn der Herr verzeucht, so hoffe und warte auf ihn, er wird kommen und nicht außen bleiben. Denn wie David sagt Ps. 147, 11.: Der Herr hat Gefallen an denen, die ihn fürchten und auf seine Güte warten. Und St. Paulus spricht Röm. 5, 5.: Die Hoffnung läßt nicht zu Schanden werden. Dies lehret auch das 3. Kapitel.

\section{Abner von Joab umgebracht}

\section{Summarium}

Die ist eine Historia, wie es sich anspinnet, daß des Sauls Nachkommen gar zu Grunde gehen, V. 9., auf daß man lerne, es verziehe sich mit der Strafe so lange es wolle, so bleibt sie doch endlich nicht außen. Dergleichen geschieht auch mit den Frommen. Denn hie siehet man, daß sich etliche Jahr mit David verzeucht, ehe er zum Regiment kommt, daß es scheint, als würde nicht allein gar nichts daraus werden,
\quellseite{461}
sondern er müßte noch auch drüber zu Boden gehen. Aber er hält fest an Gottes Verheißung; darum muß es auch endlich ins Werk kommen und gehen, wie es Gott beschlossen hatte, daß die Bösen gestraft und die Frommen gefördert werden.

\section{Summarium}

Ist eine Historia, in welcher man abermals siehet, wie Gott sein Gericht dermaßen führet, daß endlich die Unguten gestraft werde. Es scheinet aus dem 26. Capitel V. 5. des 1. Buchs Samuels, daß Abner auch dazu geholfen habe, daß David ist verfolget worden; nach Sauls Tode aber ist Abner eine Ursache gewesen, daß Israel zertrennet und der meiste Theil David zum Könige nicht wollen annehmen. Auch verkündiget sich Abner in dem, daß er mit Rizpa, Sauls Kebsweib, Hurerei treibet. Solches und anders wird zuletzt gestraft, daß er unversehens von Joab erstochen wird, V. 27. Auf daß jedermann lerne in Gottes Furcht wandeln, so er der Strafe will überhoben sein. Denn weil Gott ein gerechter Richter ist, kann die verdiente Strafe nicht außen bleiben, es verhalte und verziehe sich gleich, so lange es wolle.
\quellseite{462}

\section{Votum}

Helfe der gerechte, gütige GOtt, daß wir das Böse fliehen und das Gute thun, auf daß er die Sünde an uns nicht strafen dürfe, und daß er das Gute aus Gnaden belohne, um Jesu Christi, seines geliebten Sohnes willen, Amen.

\chapter{Das 4. Kapitel}

\section{Vorrede}

Unser HErr JESUS Christus sagt Matth. 26, 52.: Wer das Schwert nimmt, der soll durchs Schwert umkommen. Hiermit lehret er, daß alle, die das Schwert gebrauchen, und wird ihnen nicht von Gesetzen oder von der Obrigkeit gegeben, daß sie als muthwillige Todtschläger, Mörder und Aufrührer wieder sollen getödtet werden. Derwegen hüte sich jedermann für Aufruhr, und niemand gebrauche das Schwert, andern zu schaden, ohne Beruf und ohne Willen der Obrigkeit, dergleichen ohne rechtmäßige Ursachen. Denn wer über sich haut, dem fallen die Späne in die Augen, und ist nie keiner der verdienten Strafe entronnen, der sich seiner, von Gott verordneten, Obrigkeit widersetzt hat. So ist die Geduld die beste Arznei, daß man Gott die Sache und Rache befehle, und wisse, ob man dies, was in Unrecht geschieht, gleich um die Obrigkeit nicht verdienet, so hat man es doch bei Gott wohl und oft verdienet. Davon lehret auch das 4. Kapitel.

\section{Jeboseth wird ermordet, und von David gerochen}

\section{Summarium}

Hie ist eine Historia, wie die zween Hauptleute Baena und Rechob des Sauls Sohn Isboseth, ihrem Herrn, mörderisch erstechen, V. 5. 6. und meinen, sie wollen David wohl damit zu Hofe kommen, V. 8. Aber David, als einer, dem das Schwert von Gott befohlen war, läßt solchen verrätherischen Mord nicht ungerochen, sondern heißt sie auch erstechen, V. 12. Also siebet man hie abermal, daß keine Sünde ungerochen bleibt. Sonderlich aber gebet dies Exempel wider die Aufrührer, so an ihre Obrigkeit die Hand legen. Denn solches soll nicht ungerochen bleiben.

\section{Votum}

Helfe GOtt, daß wir uns unterthänig in allem christlichen Gehorsam gegen unsere Obrigkeit erzeigen, und alle aufrührerische Anschläge und dergleichen alle andere Sünde, wider GOtt und den Nächsten, fliehen und meiden, in Kraft des heiligen Geistes, Amen.
\quellseite{463}

\chapter{Das 5. Kapitel}

\section{Vorrede}

Zu Gott fliehen unter dem Kreuz ist das Beste; und da soll ein betrübt Herz, wenn es ihm übel geht, also gedenken: Siehe, jetzt habe ich das Unglück, welches mich zwar sehr ängstet und martert; aber was ist es zu rechnen gegen die vielfältigen unzähligen Wohlthaten Gottes, welche er mir sonst mein Lebenlang bis hieher erzeiget hat, und noch erzeiget? Er hat mir ja mehr Gutes gethan, denn mir alle Teufel, die ganze Welt und böse Leute Arges anthun können. Habe ich Gutes empfangen von Gott, und sollte das Böse nicht auch annehmen? Und soll also ein jeder Kreuzträger seine Seele mit Geduld fassen und Gott, bis er zu seiner und rechten Zeit hilft, aufwarten. Dieses sind auch des Königs Davids Gedanken und Geduld gewesen in seinen mancherlei Nöthen; derwegen hat ihm Gott gnädiglich und wunderlich wieder geholfen, wie auch bezeuget das 5. Kapitel.

\section{Davids andere Salbung, glücklicher Zustand und zweyfacher Sieg}

a Diese Blinden und Lahmen sind ihre Götzen gewesen, welche sie zu Trost wider David auf die Mauren setzten, als ihre Patronen, die sie schützen sollten. Wie man sieht auch mit der Heiligen Bildern thut. Als wollten sie sagen: Du kriegest mit uns nicht, sondern mit unsern Göttern; beiß dich mit ihnen, sie werden dir wohl wehren. Nicht, daß sie Blinde oder Lahme gewesen haben, sondern der Geist Gottes in denen, so hernach die Historien beschrieben haben, heißt sie also. Ps. 115, 4. 5.

b Derer heißt ein Riß oder Klag, daher diese Stätte Baal Perazim Rißmann heißt, daß die Philister da gerissen sind.
\quellseite{464}

\section{Summarium}

Ist eine Historia, wie David vom ganzen Israel zum Könige gesalbet und angenommen wird, V. 3.; auf daß man sehe, ob gleich Gott mit seiner Verheißung lange verzieht, daß er doch endlich nicht außen bleibet. Darum soll man harren und nicht müde werden.

Daß aber David so viel Weiber nimmt, V. 13., geschieht nicht aus Fürwitz, wie es die Welt deutet, sondern der Verheißung halben, daß er wußte und glaubte, Christus würde aus dem Stamme Juda kommen. Neben dem ist auch wohl zu merken, wie er alle Dinge mit Gott, und nicht aus eigenem Gutdünken anfähet; darum ist auch Glück und Heil bei ihm, daß er seinen Feinden obliegt, V. 19. So ist es wohl zu glauben, da er in solche Hoheit kommen, daß viel armer Freunde sich zu ihm gehalten und die nächsten haben wollen sein, weil die Juden ein sonder Gesetz hatten, zu den nächsten Freunden zu heirathen, wo ohne leibliche Erben der Ehemann abgestorben war.

\section{Votum}

Gott der Allmächtige erhalte uns auch in wahrem Glauben an seine Verheißungen, und helfe uns an Leib und Seele, daß wir den Unsern und andern nothleidenden Christen wieder helfen können, in leiblicher und geistlicher Noth, Amen.

\chapter{Das 6. Kapitel}

\section{Vorrede}

Gott wendet doch endlich sein gnädig Angesicht zu den Seinen, obschon der Teufel mit seinem häßlichen schrecklichen Rachen, und der Tod mit seinem grausamen Anblick, und die Welt mit feurigen, zornigen Augen uns ansehen, anmurren und wider uns toben. Derwegen wer in wahrer Buße lebt und betet, und hat Gottes gnädige Augen, der darf sich für des Teufels, des Todes und der Welt häßlichem, erschrecklichem Gesicht nicht fürchten. Denn es heißt, wie St. Petrus sagt 1. Epistel Kapitel 3. V. 12.: Die Augen des HErrn sehen auf die Gerechten, und seine Ohren auf ihr Gebet. Das ist, die gläubigen Kreuzträger siehet Gott an mit gnädigen, lachenden Augen, und höret mit leisen, offenen Ohren ihr Klagen, Seufzen und Bitten, und errettet sie endlich mit starker Hand aus aller Noth. Das lehret Davids Exempel, von dem das 6. Kapitel.

\section{Davids Freudentanz bei Abholung der Bundeslade}
\quellseite{465}

\section{Summarium}

Ist ein fürstlich Exempel, wie David sich nicht allein des Regiments, sondern auch der Kirchen und Gottesdienstes annimmt und mit allem Fleiß und Ernst fördert, V. 2., und ist so herzlich fröhlich darüber, daß er die Lade Gottes aus dem Hause Abinadab ließe holen, daß er vor großer Freude vorher tanzet, V. 15. Solches legt ihm sein Weib Michal übel aus, V. 16. Wie noch heutiges Tages viel sind, die es für eine Unehre achten, daß große Herrn sich der Lehre und des Gottesdienstes ernstlich annehmen. Denn die Welt hält solches für eine unnöthige Fahr, der man wohl könne gerathen. Aber David sagt, er begehre noch geringer zu werden für dem HErrn, denn also werde er recht zu Ehren kommen, V. 21. Und meldet der Text sonderlich, daß die hoffärtige Michal von Gott gestraft sei, daß sie hernach ihr Lebenlang kein Kind habe gehabt, V. 23. Daher gehören auch die zwei Exempel: Usa, weil er die Lade Gottes gering hielte, und mit dem Anrühren, daß er nicht Befehl hatte, unehret, schlug ihn Gott, daß er starb. Obed-Edom, der sie ehret und herberget, ward gesegnet. Alles uns zum Exempel, daß wir uns Gottes und seines Worts nicht schämen, sondern dasselbige in Ehren halten, und uns lieb lassen sein.

\section{Votum}

Wir danken GOtt, der als ein gütiger Vater das Gute belohnet und das Böse als ein rechter Richter strafet; der lehre uns durch sein Wort und regiere uns durch den heiligen Geist, daß wir alle Sünden fliehen, und allezeit Gutes wirken, Amen.

\chapter{Das 7. Capitel}

\section{Vorrede}

Ach ein jeder nehme sich der christlichen Kirche an, so viel, als er Amts, Standes und Vermögens halben kann und soll, auf daß dieselbe mit Lehre und Gottesdienst allzeit recht und christlich bestellet sei. Denn welcher Mensch ihm läßt befohlen sein der Kirchen Wohlfahrt, und fördert die mit Beisteuer, Almosen, gutem Rath, herzlichem Gebet und andern treuen Handreichungen, den wird Gott also segnen, daß er zeitlich und ewig wird dafür zu danken haben. Denn also lautet davon die wahre Verheißung unsers HErrn Jesu Christi, Matth. 10, 42.: Wer dieser Geringsten einem nur mit einem Becher kaltes Wassers tränket in eines Jüngers Namen, wahrlich, ich sage euch, es wird ihm nicht unbelohnet bleiben. So ist dies auch gewiß, ein jeder Christ kann in seinem Beruf ja etwas der Kirchen Gottes dienen und helfen. Davon erinnert auch das 7. Kapitel.
\quellseite{466}

\section[David bekommt die Verheißung des Messias, als er]{David bekommt die Verheißung des Messias, als er Willens ist, einen Tempel zu bauen}

\section{Summarium}

Das ist ein fürstlich Exempel, wie David sich mit Ernst annimmt, daß er den Gottesdienst fördert und für unbillig achtet, daß er in einem Cedern-Hause, die Lade GOttes aber unter den Teppichen wohnen soll. V. 2. Solche schöne und feine Gedanken zeigen an ein gottseliges Herz, daß er Gottes Ehre mit allem Fleiß gesucht und gerne gefördert hat. Aber Gott läßt ihm durch Nathan sagen, er wolle seinem Samen, der von seinem Lenden kommen soll, nach ihm erwecken, der seinem Namen solle ein Haus bauen, V. 2. Das ist nicht von Salomo, sondern von Christo geredet, der ihm hie verheißen wird, daß er aus seinem Samen kommen sollte, und ein ewig Reich der Gnaden anrichten. Aus diesem Text hat David den schönen herrlichen Psalm, an der Zahl der 89., V. 5. und 30. und andere mehr gemacht.

a Das ist, so redest mit mir von solchem ewigen Reich, da niemand kann König sein, er muß Gott und Mensch sein, weil er mein Sohn und doch für und für soll König sein, welches allein Gott gebühret.
\quellseite{467}

\section{Summarium}

Es ist ein sehr schön Exempel, wie David die herrliche Wohlthat erkennet und Gott dafür so herzlich danket, daß er nicht allein ihm bisher so lange viel Gutes gethan, sondern auch aufs künftige ihm Christum verheißen. Solche unaussprechliche Gnade preiset David und bekennet, daß ers nicht verdienet, sondern Gott aus lauter Gnaden um seiner Verheißung willen gethan habe, V. 21. Und bittet, V. 25., Gott wolle solcher seiner Verheißung um seiner Güte willen nachkommen. Uns zum Exempel, daß wir auch sollen lernen, erstlich erkennen, daß alles, was uns GOtt für Gutes thut, aus lauter Gnaden thue, und wir hernach Gott dafür danken, und ferner bitten sollen, daß er bei uns sein und uns weiter erhalten, und für Sünden behüten wolle.

\section{Votum}

GOtt helfe, daß die Obrigkeit auch heute seine großen väterlichen Wohlthaten erkenne, und regiere sie, daß sie die wohl und glückselig zu seiner Ehre und den Unterthanen zu Nutz allzeit recht gebrauche, hoch gelobt in alle Ewigkeit. Amen.

\chapter{Das 8. Kapitel}

\section{Vorrede}

Die Weisheit Gottes saget in Sprüchen Salomonis im 1. Capitel V. 32.: Daß die Albernen gelüstet, tödtet sie, und der Ruchlosen Glück bringet sie um. Wer aber mir gehorchet, wird sicher bleiben und genug haben, und kein Unglück fürchten. In diesen Worten wird auch folgende Lehre begriffen: Wenn ein Mensch am sichersten ist, so ist ihm allerlei Unglück und sein Untergang am nächsten. Dagegen, wenn ein Christ am schwachgläubigsten ist, und trägt das größte Kreuz, wo er nur bei Gott mit dem Gebet anhält, so ist ihm auch am nächsten seine Erlösung und Wohlfahrt. Dieses soll schrecken die Sichern und Unbußfertigen, daß sie wachen und beten, ihre Lenden umgürtet haben und ihre Lichter brennen lassen, Luc. 12, 35. Und soll trösten die Blöden und Zaghaften, und daß sie auch wachen, im Glauben stehen, männlich und stark sein, 1 Cor. 16, 13. Davon erinnert auch das 8. Kapitel.

\section{David sieget, und bestellet die Aemter}
\quellseite{468}

\section{Summarium}

Ist eine Historia, wie Gott den frommen David, nachdem er so lange im Elend gewesen und mancherlei Unglück versucht hat, gnädiglich bedenket und viel größere, herrlichere Siege gibt wider seine Feinde, V. 1. 2. Zum Exempel aller Obrigkeit, wenn sie Glück und Heil will haben, daß sie es von Gott erwarten und sich gegen Gott dermaßen halten soll, daß er nicht Ursache habe, sie zu strafen. Darum meldet der Text zweimal: Der Herr half David, wo er hinzog, V. 7. 14. Als wollte er sagen: Wo Gott mit seiner Gnade und Hülfe bei David nicht gewesen wäre, würde er mit seiner Weisheit, Sorge und Gewalt wenig ausgerichtet haben.

\section{Votum}

Der König aller Könige, unser lieber Gott, wolle heute auch aller christlichen Obrigkeit zu Schutz und Hülfe beistehen, und wolle sie regieren, daß sie sich auch ferner also halte, auf daß sie von Gott zur ewigen Herrlichkeit erhalten und erhaben werde, Amen.

\chapter{Das 9. Kapitel}

\section{Vorrede}

Wer Gott fürchtet, der stehet nicht allein wohl bei Gott für seine Person, sondern seine Kinder, Kindeskinder und Nachkommen werden auch von Gott gesegnet, das ist, sie werden für Gott groß und bei allen Engeln und Heiligen hoch geachtet, zu großen Ehren und Würden erhoben; es gibt ihnen auch Gott zeitliche Nothdurft, also, ob sie es gleich nicht haben schlechtlich, so haben sie es doch leidlich, und ob es schon nicht fleußt, so rinnet es doch, und das Wenige segnet ihnen Gott zu ihrem und der Ihren Genügen und Auskommen. Derwegen spricht David im 112. Psalm V. 2.: Das Geschlecht der Frommen wird gesegnet sein. Und sein Sohn Salomon sagt in Sprüchen im 11. Kapitel V. 21.: Der Gerechten Same wird errettet werden. Deß ist ein Exempel das Geschlecht Jonathan, davon das 9. Kapitel.

\section{Gutthätigkeit Davids gegen Mephiboseth}
\quellseite{469}

\section{Summarium}

Ist ein schön Exempel, wie David, nachdem er das Regiment zusammen gebracht und gestillet war, die Nachkommen Sauls bedenkt um Jonathans willen, V. 1., der alle Treue an ihm gethan hatte, V. 2. Denn Gottes Segen soll man niemand zu Schaden, sondern jedermann zu Nuß brauchen, auch den Feinden, wenn sie sich erkennen und Gnade begehren.

\section{Votum}

Gott verleihe uns und den Unsern auch wahre Frömmigkeit, und segne zeitlich und ewiglich alle, die an den Seinen Treue thun, und die schützen, erhalten und fördern, Amen.

\chapter{Das 10. Kapitel}

\section{Vorrede}

Wer auf Gott trauet in billigen guten Sachen, der bauet wohl und also, daß der Teufel und die Welt nicht schaden und den Bau nicht einreißen können. Derwegen ein frommer gottfürchtiger Mensch, er sei so geringe, verachtet, elend und arm in der Welt, als möglich ist, wenn er gleich auch der elendste Bettler wäre, dennoch ist er bei Gott größer und glückseliger, und auch wider den Teufel und die Welt gewaltiger und mächtiger, denn alle Gottlosen, wie groß, hoch, ansehnlich, stark und muthig sie auch sind. Derwegen wir als christliche geistliche Kriegsleute, wenn wir streiten und kämpfen, nicht nachgeben sollen, da wir nach dem Krieg wollen gekrönet werden. Sintemal Kraft und Stärke uns mittheilen will Gott, der Herr, der unaussprechliche Schätze und Güter bereitet hat im ewigen Leben denen, die da überwinden. Davon in der Offenbarung Johannis im 2. Kapitel V. 7. also gesagt wird: Wer überwindet, dem will ich zu essen geben von dem Holz des Lebens, das im Paradies Gottes ist. Und abermals V. 11.: Wer überwindet, dem soll kein Leid geschehen von dem andern Tode. Darauf höret ferner an das 10. Kapitel.

\section{Die Ammoniter verhöhnen Davids Knechte, und werden geschlagen}
\quellseite{470}

\section{Summarium}

Ist eine Historia, wie die Gottlosen, ob sie gleich bisweilen anderer Leute halben wohl möchten Friede haben, dennoch ihnen selbst Unglück und Unruhe schaffen müssen, auf daß die Sünde ja zeitlich und wohl gestraft werde. Also hätte der Ammoniter König David wohl mögen zum freundlichen Nachbarn haben; aber er und die Seinen waren's nicht werth. Darum müssen sie nach ihrem eigenen Unglück grübeln und sich zum David nöthigen, daß sie gestraft und geschlagen werden.

\section{Votum}

Steure und wehre der starke und gerechte GOtt allen, die zu Krieg und Unruhe Lust haben, und verleihe auch unserer lieben Obrigkeit Glück und Sieg wider alle ihre Feinde, und erhalte uns den heilsamen Land- und Stadtfrieden, Amen.

\chapter{Das 11. Kapitel}

\section{Vorrede}

Der heiligen Bibel fürnehmste Kapitel sollen ihnen fromme Herzen wohl bekannt machen, daß ihnen ja an nöthiger Lehre, Trost und Warnung zu keiner Zeit mangele. Ein solch Kapitel ist das 11. des andern Buchs Samuelis, welches lehret, daß wir Menschen sehr schwach und der Teufel gar stark und listig ist, der bald einen Menschen fällen kann und zu einem großen Sünder machen. Es vermahnet auch, daß niemand in seinen Sünden verzweifeln, sondern Buße thun und bei Gott Gnade suchen solle, der werde ihm Gnade erzeigen und alle seine Sünde vergeben; wie solches der Prophet Moses von Gott rühmet, 2 Mos. 34, 6., und sagt: Herr, Herr, Gott, barmherzig und gnädig, und geduldig, und von großer Gnade und Treue, der du vergießest Missethat, Uebertretung und Sünde. Wie auch dieses lehret Sirach im 18. Kapitel V. 28. und spricht: O wie ist die Barmherzigkeit des Herrn so groß, und läßt sich gnädig finden denen, so sich zu ihm bekehren. Das gemeldete Kapitel lautet, wie folgt.

\section{Davids Ehebruch und Todtschlag}
% Scan page 471: Bible verses/narrative, running heads, folios and bare cross-references only; no non-biblical apparatus retained.
\quellseite{472}

\section{Summarium}

Die ist eine sonderliche Historia, wie David sich so jämmerlich versündiget, daß er nicht allein ein Ehebrecher wird, V. 4., sondern auch bestellet, daß der unschuldige fromme Ritter Urias, da er ihn mit Güte nicht dazu vermögen könnte, daß er sein Schanddeckel gewesen wäre, erschlagen wird, und andere mit ihm, V. 17. Dadurch denn Ursache gegeben wird den Ammonitern, Gott zu lästern. Solch Exempel soll uns erstlich eine Warnung sein, daß wir nicht sicher sein, sondern stets beten sollen, daß uns Gott für Anfechtung und Sünde behüten wolle. Zum andern soll es uns auch schrecken, daß wir dergleichen Sünde nicht begehen, weil sie Gott so greulich strafet, wie hernach wird folgen.

\section{Votum}

Der treue, barmherzige Gott steure dem boshaftigen Menschenfeinde, dem Teufel, behüte uns für Sünden, richte auf, tröste und erhalte die Gefallenen, daß sie ihn mit uns loben und preisen, auch selig werden, um Jesu Christi, seines geliebten Sohns, unsers Herrn willen, Amen.

\chapter{Das 12. Kapitel}

\section{Vorrede}

Der königliche Prophet David sagt im 84. Psalm V. 5. 11.: Wie lieblich sind deine Wohnungen, Herr Zebaoth. Wohl denen, die in deinem Hause wohnen; denn ein Tag in deinen Vorhöfen ist besser, denn sonst tausend. Ja freilich ist nichts besseres, nichts herrlicheres, nichts tröstlicheres, nichts nützlicheres, denn ein wahrhaftiger Bürger und lebendiges Gliedmaß sein der rechten und wahren Kirchen Christi, da Gott wohnet und läßt sein Wort immer fort und fort erschallen, welches der rechte lebendige Brunnen und frische Quell ist, damit der Herzen oder Seelen Durst gestillet und Seele und Leib in Anfechtung erquicket werden. Dafür sollte keiner türkischer Kaiser werden, so es in seiner Willkür stünde. Sintemal eine einige rechte Predigt und Lection Gottes Worts, darin das Herz getröstet und getroffen wird, besser ist, denn aller Welt Kleinodien, Schätze, Wollust und Freude. Das hat David gelehret und recht erfahren in seinem Elende, als davon zeuget das 12. Kapitel.

\section{Nathans Bußpredigt}

\section{Summarium}

Sie ist eine Historia, wie Gott durch den Propheten Nathan David um seine begangene Missethat strafet, V. 2. Denn gleich wie Gott dem Predigtamt befohlen hat, die Sünde zu vergeben, wo mans erkennt und Gnade begehrt,
\quellseite{473}
also hat er auch befohlen, die Sünde zu strafen, wo man’s nicht erkennet. Und hie siehet man, wie greulich Gott den Ehebruch pflegt zu strafen. Sintemal der Prophet dräuet, Gott werde solche Sünde Davids mit Aufruhr strafen, V. 11.

a Jedidia) heißt lieblich dem HErrn.

\section{Summarium}

Ist endlich ein Exempel der rechtschaffenen Buße, nämlich, daß man die Sünde bekennen und um Christus willen Gnade begehren und hoffen soll. Da will Gott Gnade widerfahren lassen, wie der Prophet zu David spricht: Du wirst nicht sterben, der HErr hat deine Sünde weggenommen, V. 14.

Zum andern ist auch hie ein Exempel, daß sonderlich in die Haustafel gehöret, nämlich, daß bisweilen die Kinder der Eltern Sünde entgelten müssen. Also mußte das Kind darum sterben, daß David die Ehe gebrochen hatte, V. 15. Auf daß die Eltern lernen ihrer Kinder halben desto gottsfürchtiger und frömmer sein und sich für Sünden hüten, welche Gott auch an den Kindern bis ins dritte und vierte Glied will strafen, mit Krankheit, Armuth und andern Unglück.

\section{Votum}

Wir danken GOtt herzlich, daß wir auch das Strafgesetz und Trostevangelium haben; der verleihe Hülfe und Gnade, daß wir beide zu christlichem Leben und beständigem seligen Trost allzeit gebrauchen, bis zum seligen Todesstündlein, Amen.
\quellseite{474}

\chapter{Das 13. Kapitel}

\section{Vorrede}

Mancher meinet, es sei alles verloren, was er gutes und christliches thut, wenns ihm in diesem Leben nicht nach seinem Willen vergolten wird; mancher meinet auch, es sei alles erlogen, was Gott der Bosheit und Sünden geträuet, wenns nicht hie zeitlich, an allen nach ihrem Verdienst, gestraft wird. Aber wer mit den Augen recht siehet, mit den Ohren recht hört und mit dem Herzen, was recht ist, versteht, der wird sehen, hören und befinden, daß, was Gottes Wort lehret und desselben Dräuungen endlich erfüllet werden, als auch dieses, daß der eifrige Gott über die so ihn hassen, die Sünde der Väter heimsucht an den Kindern bis ins dritte und vierte Glied, welches er selber träuet 2 Mos. 20, 5. Dies bezeuget auch das 13. Kapitel.

\section{Der Blutschänder Amnon wird von Absaloms Dienern ermordet}

*Ledig) das ist, sie hielt sich immer, ging nicht mehr wie die Jungfrauen, im Hause unter die Leute.
\quellseite{475}

\section{Summarium}

Hier gehet die Historia an, wie dem David, da er im Regiment Friede und Ruhe hatte, seine eigenen Kinder eine Anfechtung und Unglück über das andere machen, V. 1. Auf daß wir lernen, daß solch Kreuz mit ungerathenen Kindern bisweilen auch frommen Eltern widerfähret, da es doch an Zucht und fleißigen Vermahnungen nicht gefehlet hat. Darum bedarf es nicht allein großen Fleiß, daß man die Kinder in der Zucht und Vermahnung in dem Herrn aufziehe, sondern auch stetiges Beten, daß sie wohl gerathen.

Neben dem stehet man hie, daß der Teufel durch die schändliche That, daß Ammon seine Schwester Thamar schwächte, gedenkt Mord zu stiften zwischen Brüdern, wie hernach geschah, welches David als ein Vater ein schwer Kreuz und Herzeleid ist gewesen, wie weiter folgen wird, und wir in täglicher Erfahrung auch sehen.

\section{Summarium}

Ist das andere Unglück, daß David an seinen Kindern erlebet. Nachdem Ammon seine Schwester Thamar geschwächt hatte, erwürget ihn darum sein Bruder Absalom, welches rechte Schwester die Thamar war. Sie mögen wir lernen, wie endlich Gott Hurerei strafet. Denn obwohl Absalom Unrecht thut, so hat dennoch Ammon solche Strafe wohl verdienet.

Darum soll man sich für Hurerei und Unzucht hüten, und die Eltern nicht allein die Kinder mit Fleiß davon abziehen, und allerlei Ursache, mit Gesellschaft zechen, reden, und was mehr sein kann, verhüten, sondern auch mit Ernst beten, daß Gott für solchen gefährlichen Teufels-Stricken ihre Kinder behüten wolle, da allwege Strafe und Herzeleid darauf folget.
\quellseite{476}

\section{Votum}

Gott, der HErr, regiere alle Eltern, daß sie durch Sünde ihnen und ihren Kindern Gottes Zorn und Strafe nicht auf den Hals ziehen, und erleuchte alle liebe Kinder, daß sie mit ihrem Sündenleben ihren Eltern nicht Herzeleid machen, und sich ins Verderben führen, Amen.

\chapter{Das 14. Kapitel}

\section{Vorrede}

Die Eltern strafen mit Ruthen, die Prediger mit scharfer Gesetzespredigt, die weltliche Obrigkeit aber gebraucht das Schwert, das ist, allerlei Strafe, damit sie ihren Ernst gegen die Ungehorsamen erzeigen, an derselben Ehre, Leib und Gut. Nun sagt aber St. Paulus Röm. 13, 4.: Die Obrigkeit trägt das Schwert nicht umsonst. Das ist, die Obrigkeit soll nicht allein das Schwert in ihrer Verwahrung im Kasten liegen, oder an der Wand hängen lassen, sondern sie soll es tragen, das ist, brauchen und nicht feiern lassen, und soll es tragen nicht umsonst oder vergebens, allein zur Schau oder zur Zierde, daß es also solle wohl stehen, auch nicht ohne Unterschied, sondern die Frommen damit zu schützen und die Bösen damit zu strafen. Solcher Regenten etliche hat Gott allzeit gegeben; ihm sei Lob, Ehre und Dank dafür; und ein solcher ist auch gewesen der König David, von dem das 14. Kapitel.

\section{Absalom findet auf Fürbitte des Weibes von Thekoa Gnade}
\quellseite{477}

\section{Summarium}

Sie ist ein fein Exempel, wie David einen sondern Ernst gebraucht habe wider seinen Sohn Absalom, um der Mißthat willen, daß er seinen Bruder Amnon todt geschlagen hatte, V. 1. Und scheint, wo David von dem klugen Weib von Thekoa nicht also hinterstichlichen und ihm die Zusage verdeckter Weise wäre abgedrungen worden, V. 12., daß Absalom hätte sein Lebenlang in der Fremde und Elend müssen bleiben. Solches thut David nicht allein als ein Vater, sondern auch als ein Regent. Denn denselbigen gebührt sonderlich, daß sie keine Freundschaft ansehen, sondern das Uebel sowohl an Freunden als an Fremden ernstlich strafen sollen.

\section{Summarium}

Ist wieder ein Exempel des sondern Ernstes, so David wider seinen Sohn Absalom gebraucht hat, daß, nachdem er aus Fürbitte wiederkommen ist, dennoch David ihn nicht will für sein Angesicht kommen lassen. Wo die Obrigkeit dergleichen Ernst wider die Freunde brauchte, würden die andern desto mehr Scheu haben und sich desto steifer für aller Untugend hüten. Denn es ist sehr ärgerlich und schädlich in Regimenten, wo man mehr auf Freundschaft, denn aufs Recht siehet. Denn da begibt sichs, daß man auch Fremde muß schonen und die Untugend an ihnen nicht strafen darf. So doch Gott der Obrigkeit darum das Schwerdt in die Hände gibt, daß sie es wider alles Aergerniß brauchen und wehren soll, es sei Freund oder Feind, Bekannte oder Unbekannte. Denn da soll es alles gleich gelten, nur daß Gottes Wille geschehe und dem Aergerniß mit Ernst gewehret werde.
\quellseite{478}

\section{Votum}

Gott steure allen Aufrührern und Ungehorsamen, und regiere unsere liebe Obrigkeit, daß sie ferner alles Böse ernstlich strafe, und die Frommen treulich schütze und fördere; verleihe uns allen auch, daß wir der Obrigkeit schuldigen, christlichen Gehorsam allzeit erzeigen, Amen.

\chapter{Das 15. Kapitel}

\section{Vorrede}

Beim heiligen Propheten Micha am 7. Kapitel V. 9. sagt das nothleidende israelitische Volk Gottes: Ich will des Herrn Zorn tragen, denn ich habe wider ihn gesündiget, bis er meine Sache ausführe und mir Recht schaffe. Er wird mich ans Licht bringen, daß ich meine Lust an seiner Gnade sehe. Das ist so viel gesagt: wir haben ja wohl verschuldet und verdienet die Noth, die uns betreten hat. Denn wir sind für Gott schuldig und der Strafe würdig; aber er wird uns, die wir im Finstern, aller Trübsal und Angst sitzen, wiederum aushelfen, unser Elend wenden und enden, daß uns seine Güte und Barmherzigkeit gleich als ein neu Licht wird aufgehen, und wir aus allem Herzeleid errettet, unsern Trost und Freude an seiner gnädigen, väterlichen Hülfe sehen und haben werden. In solcher Geduld, Trost und Gehorsam soll auch heute ein jeder Mensch sein Kreuz, das er mit seinen Sünden nur wohl verdient hat, tragen, als denn auch gethan hat David, von dem das 15. Kapitel.

\section{David fliehet vor seinem aufrührischen Sohn Absalom}

a Diese vierzig Jahr rechnen wir von der ersten Salbung Davids; denn er ist zwanzig Jahr alt gesalbet, und zehn Jahr Verfolgung gelitten, und im dreißigsten Jahr König beständig und angenommen.
\quellseite{479}

\section{Summarium}

Dies ist das dritte Unglück, das David von seinen Kindern widerfähret: daß Absalom, sein Sohn, einen Aufruhr wider ihn erwecket, ihn des Landes verjaget, daß er aus Jerusalem weichen und für ihn fliehen muß, V. 14. Solches Unglück ist eine Strafe des Mordes und Ehebruchs, welchen David an Uria und seinem Weibe begangen hat, wie droben der Prophet Nathan geweissaget hatte. Darum soll jedermann, sonderlich aber die Obrigkeit, sich für solchen Sünden hüten, und dies Exempel wohl merken. Denn obgleich Gott solche Sünde dem David vergeben hatte, daß sie ihm am ewigen Leben nicht schaden und er wiederum in Gottes Gnade sollte sein, darum, daß ers ihm ließe leid sein, und glaubet der Zusagung Gottes; dennoch mußte das Aergerniß, welches er damit gegeben, öffentlich gestraft und David jedermann zum Exempel werden, daß andere für dergleichen Sünde sich hüten lernten. Insbesondere wird hie sein angezeiget, wie die Aufrührer der Obrigkeit übel nachreden und sich darnach also in die Leute flechten, als wollten sie es besser machen, da gehet denn der blinde Pöbel hinan, wie wir in der Bauern Aufruhr auch erfahren haben.
\quellseite{480}

\section{Summarium}

Ist ein sonderlich Exempel, wie David erstlich die Aufruhr willig annimmt und geduldig trägt als eine Strafe, welche er mit dem Mord und Ehebruch verdienet hatte, wie ihm denn Gott durch Nathan solches dräuen ließe. Daher sehen wir, daß Gott keines Menschen verschonet, er sei niedrigen oder hohen Standes. Wer da sündiget, der bleibt gewiß nicht ungestraft. Noch ist menschliche Natur, daß wohl Gottes schreckliche Dräuung höret, daß er das Böse strafen wolle, und dies und dergleichen Exempel, der die Bibel voll ist, lieset, ja auch täglich erfährt, daß Gott das Böse straft, so verstockt und erstarrt, daß sie sich nichts daran kehret. Das ist ja ein unsäglicher Jammer.

Zum andern ist es fein, daß David in solcher Strafe und Unglück dennoch sich aller Gnade zu Gott versiehet, und es ihm heimstellet, er mache es mit ihm, wie er wolle, so soll es ihm gefallen, V. 25. Solche Geduld und Hoffnung sollen wir im Leiden und Anfechtung auch beweisen, so wir anders rechte Christen sein.

\section{Votum}

Gott sei uns gnädig und barmherzig, die wir auch unser Kreuz mit vielen großen Sünden wohl verdienet haben, und verleihe uns allen darunter Geduld und Trost, bis zur seligen von allem Uebel Erlösung, Amen.

\chapter{Das 16. Kapitel}

\section{Vorrede}

Der Prophet Jesaias sagt im 51. Kapitel V. 7.: Fürchtet euch nicht, wenn euch die Leute schmähen; und entsetzet euch nicht, wenn sie euch verjagt machen. Denn die Motten werden sie fressen, wie ein Kleid, und Würmer werden sie fressen, wie wollenes Tuch. Und im 54. Kapitel V. 14. spricht Gott selber: Du wirst ferne sein von Gewalt und Unrecht, daß du dich dafür nicht dürfest fürchten; und von Schrecken, denn es soll nicht zu dir nahen. Siehe, wer will sich wider dich rotten und dich überfallen, so sie sich ohne mich rotten? Und Lucä am 6. V. 38. sagt der HErr Christus: Eben mit dem Maß, da ihr mit messet, wird man euch wieder messen. Auch wird gesagt im Buch der Weisheit im 11. Kapitel V. 17.: Womit jemand sündiget, damit wird er auch geplagt. Dieser Warnung und Trostsprüche Exempel ist das 16. Kapitel.

\section{Ziba Bedeutung. Simei Lästerung. Bluthunde Absalom}
\quellseite{481}

\section{Summarium}

Das ist ein sonderlich Exempel einer hohen großen Geduld, daß David hören und leiden kann, wie ihn Simei so lästerlich schmähet, und will doch nicht, daß man ihm ein Leid darum thue, V. 10. Solche hohe Geduld kommet daher, daß er erkennet, wie es eitel verdiente Strafe sei. Wie ihm nun seine Sünde von Herzen leid ist, leidet er gern und willig, was ihm der Sünden halben begegnet. Solche Geduld ist auch ein Stück, das zur rechten christlichen Buße gehöret.

\section{Summarium}

Ist eine Historia, in welcher wir sehen, wie dennoch Gott allezeit Mittel schaffet in der Anfechtung, auf daß sich die Seinen nicht beklagen dürfen, sie seien ganz verlassen. Darum schaffet Gott, daß Husai, der Arachiter, aus Anregung Davids gen Jerusalem kommet, und stellet sich, als halte ers mit Absalom, und das darum, daß er Ahitophels (der auf Absaloms Seiten war) Rathschlag zunichte machen und David warnen könnte.

Zum andern siehet man hie ein sonderlich Exempel, daß über das, daß David seines Ehebruchs und Todschlags halben des Landes verjagt wird, Gott auch denselben Ehebruch damit strafet, daß gleich wie er einem andern sein Weib heimlich genommen und sie beschlafen hatte: also auch sein Nächster (das ist Absalom) bei seinen Kebsweibern schläft öffentlich an der lichten Sonne, 2 Sam. 12, 11. Nach dem Spruch Weisheit. 11. V. 17.: Damit jemand sündiget, muß er auch gestraft werden. Auf daß jedermann sehe, wie Gott allerlei Unzucht und sonderlich dem Ehebruch feind ist, und sich desto fleißiger dafür hüte.

\section{Votum}

Der treue, barmherzige GOtt verleihe uns auch, daß wir in Geduld alle Lästerung überwinden, offenbare auch alle schädliche Rathschläge, und helfe, daß wir uns für Unzucht und andern Sünden hüten, damit wir nicht verdiente Strafe leiden dürfen, Amen.

\chapter{Das 17. Kapitel}

\section{Vorrede}

List und Gewalt ist der Feinde Gottes Rüstung, damit sie wider seine Kirche streiten. Dagegen Beten und geduldig sein, sind der Christen Waffen, damit sie wider Gottes und ihre Feinde kämpfen. Und Gott sei Lob in Ewigkeit, dieser Theil gewinnet, jener aber verlieret, obschon die Gottlosen eine Weile triumphiren, und die Kinder Gottes eine Zeitlang leiden müssen. Derwegen von beiden also sagt David im 92. Psalm V. 8.: Die Gottlosen grünen wie das Gras, und die Uebelthäter blühen alle, bis sie vertilget werden, immer und ewiglich. Und V. 13.: Der Gerechte wird grünen wie ein Palmbaum, er wird wachsen wie eine Ceder auf Libanon. Davon höret ferner an ein Exempel aus dem 17. Kapitel.

\section{Ahitophels nichtiger Rath, und verzweifeltes Ende}
\quellseite{482}

\section{Summarium}

Ist eine Historia, in welcher man siehet, mit was Ernst und großem Fleiß der Teufel dem frommen David nachstellet daß (menschlicher Weise davon zu reden) nicht möglich ist, daß er entgehen könnte. Aber Gott schaffet dennoch Mittel, daß des Teufels und böser Leute geschwinde Anschläge müssen zu nichte werden. Wie man hie an Ahitophel siehet, der einen Rath gab, der Absalom und allen Aeltesten in Israel gut däuchte; noch macht Gott durch Husai ihn zu nichten, V. 7. Darum liegt es an dem, daß jedermann auf Gott und sein Wort sehe, ihn mit Ernst anrufe und ihm alle Sache befehle; wie David hie thut, da er in dieser seiner großen Noth also zu Gott schreiet: Herr mache den Rathschlag Ahitophels zur Narrheit, 2 Sam. 15, 13., Ps. 31, 18.; so hilft denn Gott gewißlich.
\quellseite{483}

\section{Summarium}

Ist eine sonderliche Historia, wie es mit dem weisen Mann Ahitophel ein Ende nimmt, welcher ein Exempel ist aller deren, so mit menschlicher Weisheit hoch für andern begabt sind, aber doch derselbigen mißbrauchen zur Tyrannei wider das Wort Gottes und die christliche Kirche. Demselben soll es also gehen, wie Salomon sagt, Sprüche 21, 30.: Es hilft keine Weisheit, kein Verstand, keine Kunst wider den Herrn. Darum, weil sie nicht ablassen wollen, müssen sie darüber zu Boden gehen. Wie denn dem Ahitophel sein weiser Rathschlag Ursache gibt, daß er hingehet, und für Leid sich selbst erkennt, daß derselbige nicht fortgegangen war.

\section{Votum}

Gott, der Herr, mache auch heute zunichte des Teufels und aller Gottlosen böse und schädliche Anschläge, und rette und erhalte seine Christenheit und Gläubige in aller Gefahr und Noth, um Jesu Christi, seines Sohnes willen, Amen.

\chapter{Das 18. Kapitel}

\section{Vorrede}

Ungehorsame Kinder laden auf sich diese schwere Bürde und Sünden: Gottes und der Eltern Segen wird von ihnen, ihr Glück aber über sie sein, Sir. 3, 9. 10. Mit bösem Gewissen müssen sie sich tragen ihr Lebenlang; im Sterbestündlein wird ihnen ihr Licht verlöschen mitten im Finsterniß, Spr. Sal. 20, 20. Aus dem Himmel sollen sie verstoßen sein ewiglich, es sei denn, daß sie Buße thun; in der Hölle müssen sie bleiben, in Qual und Pein, in alle Ewigkeit. Das ist eins. Darnach, kein Schmerzen ist so groß, nach dem Schmerzen, denn ein Herz fühlet, wenn dasselbe Gottes Zorn, wegen begangener Sünde, brennet und drücket, als der Schmerzen ist, den die Eltern empfangen, haben und fühlen, wenn es ihren Kindern übel gehet, oder wenn sie in Noth stecken und übel umkommen, sonderlich wo man auch der Seelen Verlust muß fürchten. Beides lehret auch das 18. Kapitel.

Absalom wird an einer Eiche hangend durchstochen.
\quellseite{484}

a Das ist, Absicht.

b Vernimm, Leide thue.

\section{Summarium}

Dies ist eine Historia, wie es mit dem Aufrührer Absalom und denen, so ihm zur Aufruhr geholfen, gegangen und ein Ende genommen habe, nämlich, daß in die zwanzig tausend Mann umgekommen, V. 7., und Absalom mit dem Kopf, durch ein wunderbarlich Urtheil Gottes, an einer Eiche hangen bleibt und also dran erstochen wird, V. 9. Neben solcher Strafe zeiget der Text noch eins an, wie Absalom, der seinen Vater so übel gehalten hat, nicht werth ist gewesen, daß er einen Sohn hat sollen haben, und ist also ohne Erben ausgerottet, V. 14. Darum sollen erstlich die Kinder lernen, daß sie ihre Eltern ehren und für Augen haben. Zum andern sollen alle Unterthanen lernen sich in Gehorsam halten und für Aufruhr zu hüten, als die nimmermehr ohne sonderlichen großen Schaden Land und Leute pflegt abzugehen.
\quellseite{485}

\section{Summarium}

Ist eine Historia, wie David um seinen Sohn Absalom sehr trauert, der ihn doch erstlich hoch bekümmert, darnach des Landes verjagt und endlich ihm nach Leib und Leben gestanden hatte. Aus welchem allem man leichtlich kann abnehmen, wie er solche Sünde alle von Herzen gerne vergeben und ihn wieder zu Gnaden wollte angenommen haben. Nicht allein deshalb, daß er sein Sohn war, sondern auch darum (wie vorgemeldet), daß David solchen Aufruhr nicht anders annahm, als eine verdiente Strafe seiner Sünden. Solche Erkenntniß der Sünden macht ein freundliches Herz, das gern und willig vergibt, auch seinen ärgsten Feinden.

\section{Votum}

GOtt, der heilige Geist, regiere die liebe Jugend, daß sie ihre Eltern nicht betrüben, sondern die ehren und ihnen gehorsam sein; tröste auch und stärke alle betrübte Eltern und bekümmerte Herzen, daß sie ihr Creuz in Geduld tragen und überwinden, Amen.

\chapter{Das 19. Kapitel}

\section{Predigt}

Welcher Mensch seinem Fleisch allen Muthwillen gestattet, und in Sünden lebet, der gehört Christum nicht an. Wer aber seinen Glauben beweiset mit christlichem Leben, nach seinem Beruf, Amt und Stande, der hat Gemeinschaft an GOtt. Derwegen sollen wir auf GOttes Gnade und Barmherzigkeit nicht muthwillig sündigen, sondern sollen gottselig leben und uns aller guten Werke befleißen. Dazu soll uns auch bewegen, daß Gott ist ein eifriger Gott, der die Sünde hasset und strafet, wie er dagegen das Gute liebet und belohnet. Es soll uns auch dazu treiben, daß solches Gott zu Lob, Ehre und Preis dienet; und weil wir fürnehmlich dazu erschaffen, darum erlöst und mit Gottes Wort alle Tage unterrichtet werden, daß wir in Gerechtigkeit und Heiligkeit leben sollen. So betrachten wir allhie auch billig die Vermahnung St. Pauli an die Philipper, da er im 4. Kapitel V. 8. also saget: Was wahrhaftig ist, was ehrbar ist, was gerecht, was keusch, was lieblich, was wohl lautet, ist etwa eine Tugend, ist etwa ein Lob, dem denket nach. So wird der HErr des Friedens mit euch sein. Daher gehört auch das 19. Kapitel.

\section{David wird wieder ins Königreich eingesetzt}
\quellseite{486}

\section{Summarium}

Sie ist eine feine Historia, wie David, da ihm Gott wieder ins Königreich hilft, alle Gnade und Barmherzigkeit jedermann erzeigt, wer Gnade begehrt und nicht tyrannisch mit den Leuten verfährt, die es doch wohl verdienet hatten, als mit Simei, V. 16., Amasa und andern. Denn dies ist die rechte Krone aller fürstlichen Tugenden, daß man Gnade beweise und nicht tyrannisch fahre, wo die, so gesündiget, ihre Sünde erkennen und Gnade begehren.
\quellseite{487}

\section{Summarium}

Dies ist ein Exempel, das sonderlich gen Hof gehöret, wie Ziba seinen Herrn Mephiboseth fälschlich bei David angegeben und in Ungnade gebracht hat, auf daß er zu Hofe desto angenehmer wäre und mehr Gunst hätte. Solches beklaget sich Mephiboseth, V. 26. Aber David beharrt auf dem, daß er Ziba zugesagt hatte, V. 29. Diese Plage ist gemein zu Hofe, daß auch fromme Fürsten sich nicht allweg für Heuchlern hüten können, und bisweilen den Leuten dadurch Unrecht geschiehet, allein darum, daß die Herren zu bald glauben, wie man auch an David siehet. Aber in dem gehet Gottes Urtheil und Strafe wider Saul, daß alle seine Nachkommen solcher Sünden entgelten und weder Glück und Segen müssen haben. Darum die Eltern auch ihrer Kinder halben sich in Gottesfurcht halten und für Sünden hüten sollen, daß die Kinder derselbigen nicht entgelten müssen.

\section{Votum}

Gott, der HErr, regiere Obrigkeit und Unterthanen; die Obrigkeit, daß sie recht und christlich auch regiere; die Unterthanen, daß sie fromm und gehorsam sein, und also auch ihrer Kinder verschonen, daß dieselben ihre Sünde, mit ihrem Schaden und Herzleid, nicht büßen dürfen, Amen.
\quellseite{488}

\chapter{Das 20. Kapitel}

\section{Vorrede}

Wie keine Sünde allein ist, es folget immer eine der andern, wo man nicht in Zeiten Buße thut: also, wenn Gott anfänget, die Leute im Zorn heimzusuchen, so folget auch eine Strafe auf die andere; und wer dem vormals in dem Sündenleben fröhlich gewesen, ach, bei dem wird das Lachen sehr theuer, unter den vielen und großen Strafen. Derwegen wollen wir durch Hülfe des heiligen Geistes aufhören, Böses zu thun und von Sünden abstehen, daß Gott nicht strafen dürfe, und wollen Gutes thun und in allen Tugenden uns fleißig üben, daß Gott uns segne und an Leib und Seele behüte aus Gnaden, um Jesu Christi, seines geliebten Sohnes willen, Amen. Dies soll auch erinnern das 20. Kapitel.

\section{Seba, der Aufrührer, kriegt seinen Lohn}

a Berühmter) einer von den großen Hansen, vom hohen Adel, der großen Anhang im Volk und Ansehen oder Namen hatte, wie Catilina zu Rom.
\quellseite{489}

\section{Summarium}

Sie ist eine Historia, wie ein anderer Aufruhr wider David erreget wird, daß allein der Stamm Juda an ihm bleibt hangen, v. 2. Auf daß man lerne Gott fürchten und sich für Sünden hüten. Denn wenn Gott einmal anfähet zu strafen, hält er lange an. Es ist auch hier ein Exempel, wie es dem Aufrührer Seba gehet, nämlich, daß er den Kopf drüber lassen muß, v. 22. Denn Gott hat’s so beschlossen: Wer das Schwert nimmt, dem es nicht gehöret, soll durchs Schwert wieder umkommen.

\section{Votum}

Der allmächtige Gott steure heute auch allen Aufrührern und ungehorsamen Unterthanen, und regiere uns, daß wir ja nicht in Sünden leben, sterben und verderben; sondern helfe, daß wir in Gottesfurcht beständig bis zum seligen Simeonis-Stündlein verharren, Amen.

\chapter{Das 21. Kapitel}

\section{Vorrede}

Viele Eltern leider! erben auf ihre Kinder nicht Gottes Segen, sondern seinen Fluch und verlassen seinen Zorn und Strafe zum Gedächtniß. Denn wie die Kinder wohl über langer der Eltern Gottes-Furcht und Gehorsam gegen Gott genießen können: also müssen auch der Gottlosen Kinder ihrer Eltern Mißthat ausfeucheln, wenn es gleich noch so lange anstünde. Dieses sollen alle Väter und Mütter wohl wahrnehmen, und in allem schuldigen Gehorsam gegen Gott wandeln; und also werden sie die besten Schätze ihren Nachkommen sammeln. Dagegen sollen sie meiden alle Sünde und Mißthat. Denn sonst führen sie also Gottes Fluch und verderblichen Schaden auf dieselben. Ferner sagt David im 34. Psalm v. 20. Der Gerechte, das ist, ein gläubiger und frommer Mensch muß viel leiden; aber der HErr hilft ihm aus dem allen. Vernehmt, durch Trost, Errettung und endliche Erlösung. Beides der Lehre und Trostspruchs Exempel ist das 21. Kapitel.

\section[David wehret der Theurung; erhält Siege wider die]{David wehret der Theurung; erhält Siege wider die Philister}
\quellseite{490}

\section{Summarium}

Ist eine Historia, wie Gott nach dem Aufruhr das Land mit Hunger gestraft habe drei Jahre lang, V. 1.

* Sonderlich aber ist zu merken, wie sieben Söhne vom Geschlechte Saul gehenkt sind worden, darum, daß ihr Vater die Gibeoniter ohne alle Schuld, wider den Eid, so ihnen von dem Volk Israel geschehen, getödtet hat, V. 8. Auf daß die Eltern auch desto gottesfürchtiger und frömmer lernen sein, daß ihre Kinder ihrer Sünde nicht entgelten und desto mehr Unglück müssen leiden, wie man sehet.

\section{Summarium}

Ist eine historia, wie David weiter mit Kriegen wird angefochten und sonderlich in einer Schlacht Noth leidet, daß, wo Abisal ihm nicht geholfen hätte, er umgekommen wäre, V. 15. Auf daß man sehe, wie der fromme König so in mancherlei Gefahr und Unglück gesteckt ist, aber dennoch Gott in allen Wegen ihm gnädiglich heraus geholfen hat. Denn, ein König oder Regent sein, heißt nicht gute Tage haben, sondern seines Berufs fleißig warten, und darüber sich keiner Gefahr weigern, die sich im Handeln täglich mit Haufen finden, wie man am frommen David siehet.

\section{Votum}

Der treue, liebe GOtt regiere die Eltern, daß sie also leben, handeln und wandeln, damit
\quellseite{491}
 sie ja ihren Kindern nicht schädlich sein; und erlöse alle, die das Kreuz in wahrem Glauben und Geduld tragen, nach seinem Willen und wie es ihnen gut und selig sein soll, Amen.

\chapter{Das 22. Kapitel}

\section{Vorrede}

In der Noth sollen wir Gott um Hülfe anrufen und wenn uns geholfen ist, sollen wir auch Gott danken für erzeigte Wohlthat. Denn also befiehlet es Gott selber im 50. Psalm V. 15.: Rufe mich an, spricht er, in der Noth, so will ich dich erretten, so sollt du mich preisen. Und 1 Tim. 2 V. 1. vermahnet St. Paulus auch dazu und sagt: So ermahne ich nun, daß man für allen Dingen zuerst thue Bitte, Gebet, Fürbitte und Danksagung für alle Menschen. Für die Könige und für alle Obrigkeit, auf daß wir ein ruhiges und stilles Leben führen mögen in aller Gottseligkeit und Ehrbarkeit. Denn solches ist gut, dazu auch angenehm vor Gott, unserm Heilande. Wie solches der heilige David verstanden und gelehret, also hat er auch selber uns zum Exempel gethan. Wie bezeuget das 22. Kapitel.

\section{Lobgesang Davids für die Errettung von seinen Feinden}
\quellseite{492}

a (Sich zu lieben) das sind diejenigen, die sich mit vielen Gottesdiensten wollen um GOtt wohl verdienen, meinen herzlich und thun mit Ernst, aber ohne Gottes Wort, aus eigen erwähltem Fürnehmen, wie unsere Mönche und alle Abgöttische thun.

\section{Summarium}

Das ist ein Dank-Psalm, darin David Gott danket, daß er ihn von allen seinen Feinden erlöset und in allerlei Noth und Gefahr gnädiglich erhalten hat, wider Saul, wider Absalom, wider die Aufrührer und wider die Heiden, welche allzumal, wie ein großes Wetter, David gern zu Boden geschlagen und ersäuft hätten. Aber weil er sich an Gottes Hülfe hält, müssen sie ihn bleiben lassen, wie er spricht: Wenn mir angst ist, so rufe ich den HErrn an, und schreie zu meinem Gott, so erhöret er meine Stimme von seinem heiligen Tempel.

Solchem Exempel sollen wir auch folgen und wissen, wer sich in der Noth an etwas anders, denn an Gottes Hülfe und das Gebet hält, muß endlich zu Boden gehen. Wohl aber denen, die auf den HErrn hoffen und seiner Güte warten. Denn ob sie wohl Unglück nicht können überhoben sein, so ist doch Gott bei ihnen, der hilft in der Noth.

\section{Votum}

Wir danken GOtt auch herzlich, der uns bisher in mancher großer Noth beigestanden und erhalten hat; der wolle ferner unser treuer GOtt und lieber Vater sein und bleiben, daß wir aus aller Noth errettet und im Tode zum ewigen freudenreichen Leben erhalten werden, Amen.

\chapter{Das 23. Kapitel}

\section{Vorrede}

Das 23. Kapitel des andern Buchs Samuels ist fürnehmlich des Königs Davids Testament und letzter Wille. Denn obwohl David hernach eine Zeitlang gelebet und viel geredet, gethan und gelitten hat, dennoch hat er dies sein Testament sein und bleiben lassen. Er lehret aber darin von den fürnehmsten Artikeln christlicher Lehre, als vom wahren Erkenntniß Gottes, von beiden Naturen und großen Wohlthaten des verheißenen Messias unsers HErrn Jesu Christi, und von wahrer Anrufung Gottes. Welches so heilsame und nöthige Artikel sind, daß der HErr Christus Joh. 17, V. 3. also davon sagt: Das ist das ewige Leben, daß sie dich, der du allein wahrer Gott bist, und den du gesandt hast, Jesum Christum, erkennen. Derwegen dies Testament Davids jedermann wohl bekannt sein solle. Solches aber lautet aus dem 23. Kapitel.
\quellseite{493}

\section{Davids letzte Reden und Heldenbuch}

a Ohne Wolken) Weiß richtet des Gesetzes Reich an auf dem Berge Sinai mit Donnern, Wolken, Blitzen schrecklich. Aber dies Reich wird herrlich sein, wie es ist im Lenz, wenn es geregnet hat und die Sonne früh scheint.

b Nichts wächst) kein Königreich ist so hoch für Gott, wird auch nicht so wachsen, sondern vergehen. Allein dies Reich bestehet ewiglich.

c Belial) sind die, so dem Reich Christi feind sind, als Juden, Papst, Ketzer, Türken etc. Die wollen allein nütze und die besten sein, und sind doch die schädlichsten, darum heißen sie Belial, die Unnützen oder Schädlichen. Also sagt Jer. 23, 3. von den falschen Propheten: Sie sind mit ihrem Nützen kein Nutzen diesem Volk, das ist, sie sind die schädlichsten, eben da sie nütze sein wollen.

d Jasabeam) an diesem Ort steht es im Hebräischen also: dies sind die Namen der Helden David, Josebeth, Tachmoni der fürnehmste unter dreien. *Jofelino, daßnitz, und schlug achthundert auf einmal. Da achten wir, der Text sei durch einen Schreiber verderbet, etwa aus einem Buch unkenntlicher Schrift, und von unsern Buchstaben. Und sei also Abino für Oret, und Hachnitz für Ethanito gemacht. Denn die Hebräer wohl wissen, wie man in ihrer Handschrift kann Daleth für Resch, Nun für Zain, He für Thau und wiederum lesen. Darum haben wir nach dem Text 1 Chron. 12, 11. corrigiret. Denn der Text an diesem Ort nichts gibt. Dergleichen kann auch geschehen sein in dem Wörtlein drei, dem acht hundert, so in der Chronica dreißig Item dreihundert stehen. Doch kann das eine andere Meinung haben, wie unten 1 Chron. 12, 11.
\quellseite{494}

\section{Summarium}

Ist ersichtlich eine Weissagung Davids von seinem Königreich, daß es aufgehen, wachsen und sein werde, wie das Licht ist am Morgen, wenn die Sonne aufgehet frühe ohne Wolken, dergleichen Reich in der ganzen Welt nicht sein werde, V. 4. Solches gehet auf das Reich Christi, welcher vom Stamme David ist herkommen. Desselbigen Feinde sollen wie die Dornen ausgerottet werden. Zum andern werden hie erzählet die tapfern Helden, welche David bei ihm zu Hof gehabt hat, V. 8. Da sehet man, wie Gott erstlich viel seiner Leute gibt und einem Fürsten an die Seite setzet, wenn er einem Reich helfen und es groß will machen. Wiederum, wo es an seinen redlichen Leuten fehlet, da werden die Regimente nicht lange bestehen können.

\section{Votum}

Gott, der HErr, segne und erhalte ferner das Reich seines lieben Sohnes, unsers HErrn Jesu Christi, und gebe auch im weltlichen Regiment seiner Leute, die zu seinen Ehren und gemeinem Nutzen zum Besten herrschen und dienen, Amen.

\chapter{Das 24. Kapitel}

\section{Vorrede}

Dieweil auch nach der Wiedergeburt in allen Menschen viel übrige Sünden im Fleisch bleiben, durch welche ein Christ, wo er nicht wohl Achtung auf sich selbst gibt, in verderbliche Sicherheit kann verführet werden; als soll derhalben ein jeder desto fleißiger auf sich selbst Achtung geben, damit er durch solche anklebende sündliche Gebrechen und Zuneigungen mit diesem Laster nicht eingenommen werde, welches denn ohne Verleugnung des Gewissens, ohne Verlierung des heil. Geistes und des Glaubens, auch ohne zeitliche und ewige Strafe nicht abgehen würde. Derwegen zur Vermeidung solches Schadens soll ein jedes durch tägliche Reue und Leid, und durch ernstlichen heftigen Widerstand, neben inbrünstigem Rufen und Schreien zu Gott, dahin trachten, damit es solchen bösen Lüsten, durch Beistand des heil. Geistes möge widerstehen, auf daß solche Funken nicht zu einem schädlichen Feuer werden, das ist, es soll ein jedes, wie St. Paulus sagt Gal. 5, 24.: Sein Fleisch sammt den Lüsten und Begierden kreuzigen. Das erinnert auch das 24. und letzte Kapitel des andern Buchs Samuels.

\section[Davids Ehrgeiz in Zählung des Volks wird mit]{Davids Ehrgeiz in Zählung des Volks wird mit der Pestilenz bestrafet}
\quellseite{495}

\section{Summarium}

Das ist wiederum ein Exempel, wie auch heilige Leute, wenn sie außer der Noth und Anfechtung sind, leichtlich in Sicherheit und Sünde fallen. Derhalben ist nütze und gut, daß Gott durchs Kreuz uns in seiner Furcht erhalte. Droben haben wir oft gehört, wie David, wenn er in Anfechtung ist, vom Saul verfolget wird oder sonst mit ihm übel zugehet, betet und geduldig ist. Wie aber, nachdem er Friede hat und sich für niemand besorgen darf, wird er sicher, will wissen, wie viel er Volks in das Feld wider die Feinde vermöge, V. 2. So doch Gott verboten hat, 2 Mos. 39, 12., man sollte das Volk nicht zählen, wollte mans aber zählen, so sollte ein jeglicher Haupt sich mit einem Opfer lösen; sonst würde ihnen eine Plage widerfahren. Solches bedenkt Joab und die andern Hauptleute, und widerrathen es, V. 3. Aber David will sich nicht bereden lassen, er gedenkt, er sei Herr, andere sollen ihm folgen, er dürfe niemand folgen. Wie denn große Herren diese Anfechtung haben und wollen ihres Sinnes sein. Es geräth aber nicht allezeit wohl, wie man an David auch siehet.
